\documentclass{usydthesis}
\newcommand{\memyself}{Chenuka Garusinghe\xspace}
% Configuration

\def\degree{Bachelor of Advanced Computing (Honours)}
\def\department{School of Computer Science}

\title{{\bf\Huge When Demonstrations Lose to Priors}\\[0.6em]
       {\Large Demonstration Selection for In-Context Learning under Distribution Shift in 7B Language Models}}
\author{\memyself}
\def\sid{000000000}          % TODO: student ID
\def\supervisor{Dr Sichao Li}   % TODO: supervisor
%\def\assocsupervisor{Comment this line to remove assoc supervisor}

%%%%%%%%%%%%
% Packages
\usepackage{alltt}
\usepackage{amsfonts}
\usepackage{amsthm}
\usepackage[colorlinks=false,plainpages=false,pdfborder={0 0 0},backref=false]{hyperref}
\usepackage[all]{hypcap}
\usepackage{listings}
\usepackage{longtable}
\usepackage{multirow}
\usepackage{pdflscape}
\usepackage{pifont}
\usepackage{rotating}
\usepackage{setspace}
% References: APA 7th edition (biblatex-apa, biber). natbib=true keeps \citep/\citet.
\usepackage[australian]{babel}
\usepackage{csquotes}
\usepackage[style=apa,backend=biber,natbib=true,uniquename=false,uniquelist=false]{biblatex}
\addbibresource{references.bib}
\newcommand{\newcite}[1]{\citet{#1}}
\usepackage{tabularx}
\usepackage{wrapfig}
\usepackage{xcolor}
\usepackage{xspace}
\usepackage{booktabs}
\usepackage{enumitem}
\usepackage{caption}
\usepackage{soul}
\usepackage{array}
\usepackage{tikz}
\usetikzlibrary{arrows.meta,positioning,shapes.geometric,calc,fit,backgrounds,decorations.pathreplacing}

% Macros and environments used in the design appendix
\newcommand{\RR}{\mathbb{R}}
\newcommand{\1}{\mathbf{1}}
\newcommand{\E}{\mathbb{E}}
\DeclareMathOperator{\sign}{sign}
\DeclareMathOperator{\sigmoid}{\sigma}
\DeclareMathOperator{\logsumexp}{logsumexp}
\DeclareMathOperator{\Maj}{Maj}
\DeclareMathOperator{\corr}{corr}
\DeclareMathOperator{\Cov}{Cov}
\DeclareMathOperator{\Var}{Var}
\DeclareMathOperator{\RMSNorm}{RMSNorm}
\theoremstyle{plain}
\newtheorem{theorem}{Theorem}[chapter]
\newtheorem{proposition}[theorem]{Proposition}
\let\lemma\relax\let\endlemma\relax
\newtheorem{lemma}[theorem]{Lemma}
\newtheorem{corollary}[theorem]{Corollary}
\theoremstyle{definition}
\newtheorem{definition}[theorem]{Definition}
\theoremstyle{remark}
\newtheorem{remark}[theorem]{Remark}
\newcolumntype{L}[1]{>{\raggedright\arraybackslash}p{#1}}
\lstset{basicstyle=\ttfamily\small, breaklines=true, breakatwhitespace=true, frame=single,
        backgroundcolor=\color{gray!5}, columns=fullflexible}
\setlength{\emergencystretch}{3em}

\def\subsectionautorefname{Section}
\def\chapterautorefname{Chapter}
\def\partautorefname{Part}

\newcommand{\tick}{\ding{51}}
\newcommand{\cross}{\ding{53}}

\newcounter{todocounter}
\newcommand{\todo}[1]{\marginpar{$|$}\textcolor{red}{\stepcounter{todocounter}\footnote[\thetodocounter]{\textcolor{red}{\textbf{TODO }}\textit{#1}}}}

%%%%%%%%%%%%
% Thesis macros (see introduction/introduction.tex for how findings are used)
% A finding: one bolded sentence, numbered by chapter, stated for the conditions tested.
\newcounter{finding}[chapter]
\renewcommand{\thefinding}{\thechapter.\arabic{finding}}
\newcommand{\finding}[1]{\par\smallskip\noindent\refstepcounter{finding}%
  \textbf{Finding \thefinding.\ #1}\par\smallskip}
% Anticipated criticism: a question an examiner may ask, and the answer.
\newcommand{\critique}[2]{\par\smallskip\noindent\textit{``#1''}\par\noindent #2\par\smallskip}
\newcommand{\AUROC}{\textsc{auroc}\xspace}
\newcommand{\RQ}[1]{\textbf{RQ#1}\xspace}
\newcommand{\sata}{\textsc{sata}\xspace}
\newcommand{\satami}{\textsc{sata-mi}\xspace}
\newcommand{\strat}[1]{\texttt{#1}}
\newcommand{\figpath}{figures}
% Parts are headings in the table of contents only (no separate title page)
\newcommand{\thesispart}[1]{\cleardoublepage\phantomsection\addcontentsline{toc}{part}{#1}}

%%%%%%%%%%%%
% Setup

%   Page size:
\oddsidemargin=0cm	% really 1in
\evensidemargin=0cm
\textwidth=6.2677165in

% initial page numbers:  i, ii, iii, ...
\renewcommand{\thepage}{\roman{page}}

\newcommand{\defn}[1]{\textit{#1}}
\newcommand{\ttl}[1]{\textit{#1}}
\newcommand{\txt}[1]{\textsf{\smaller #1}}
\newcommand{\qu}[1]{\textsf{#1}}
\newcommand{\txtbf}[1]{\textsf{\textbf{\small #1}}}
\newcommand{\quot}[1]{\textit{#1}}

\newcommand{\sssection}[1]{\textsf{\textbf{#1}}}

\newcommand{\cf}[1]{\mbox{$\it{#1}$}}   % category font

\newcommand{\alta}{\textsc{alta}\xspace}
\newcommand{\api}{\textsc{api}\xspace}
\newcommand{\candc}{C\&C\xspace}
\newcommand{\ccg}{\textsc{ccg}\xspace}
\newcommand{\ccgbank}{CCGbank\xspace}
\newcommand{\cky}{\textsc{cky}\xspace}
\newcommand{\jhu}{\textsc{jhu}\xspace}
\newcommand{\lfg}{\textsc{lfg}\xspace}
\newcommand{\hpsg}{\textsc{hpsg}\xspace}
\newcommand{\mwe}{\textsc{mwe}\xspace}
\newcommand{\mwes}{\mwe{}s\xspace}
\newcommand{\NE}{\textsc{ne}\xspace}
\newcommand{\Naive}{Na\"{i}ve\xspace}
\newcommand{\naive}{na\"{i}ve\xspace}
\newcommand{\ngram}{$n$-gram\xspace}
\newcommand{\ngrams}{\ngram{}s\xspace}
\newcommand{\nlp}{\textsc{nlp}\xspace}
\newcommand{\np}{\textsc{np}\xspace}
\newcommand{\nps}{\np{}s\xspace}
\newcommand{\parseval}{\textsc{parseval}\xspace}
\newcommand{\pos}{\textsc{pos}\xspace}
\newcommand{\pp}{\textsc{pp}\xspace}
\newcommand{\qa}{\textsc{qa}\xspace}
\newcommand{\ram}{\textsc{ram}\xspace}
\newcommand{\rasp}{\textsc{rasp}\xspace}
\newcommand{\tbl}{\textsc{tbl}\xspace}
\newcommand{\ptb}{\textsc{ptb}\xspace}
\newcommand{\wsj}{\textsc{wsj}\xspace}

\newtheorem*{assumption}{Assumption}
\newtheorem*{mainassumption}{Main Assumption}
\newtheorem*{corollary*}{Corollary}




%%%%%%%%%%%%
% Start

\begin{document}

%%%%%%%%%%%%
% Title page
\maketitle

\setstretch{1.5}

% intro pages
\cleardoublepage
\phantomsection
\chapter*{Student Plagiarism: Compliance Statement} \label{sec:plagiarism}
{
\setlength{\parindent}{0cm}
\setlength{\parskip}{1em}

I \memyself certify that:   

I have read and understood the University of Sydney Student Plagiarism:  Coursework Policy and Procedure;

I understand that failure to comply with the Student Plagiarism: Coursework Policy and Procedure can lead to the University commencing proceedings against  me for potential student misconduct under Chapter 8 of the University of Sydney  By-Law 1999 (as amended);  

This work is substantially my own, and to the extent that any part of this work is not my own I have indicated that it is not my own by acknowledging the source of that part or those parts of the work.
}

\vspace{4cm}
\begin{tabular}{llll}
\textbf{Name}: & \authors & &  \\[2em]
\textbf{Signature}: & \includegraphics[width=3cm]{images/signature.jpg} \hspace{0.4\textwidth} & \textbf{Date}: & \hspace{0.4\textwidth}\\[2cm]

\end{tabular}




\cleardoublepage
\phantomsection
\chapter*{Abstract}

Large language models can learn a task from a handful of labelled examples placed in their prompt, a behaviour called in-context learning.
Which examples are shown matters, and a growing literature selects them by similarity, diversity or learned scores.
Little is known about whether selection helps when the test data shift away from the examples, or when what the examples teach contradicts what the model already believes.
Smaller models are known to follow such contradicting evidence only partially.

This thesis asks whether choosing the demonstrations can help a 7B model follow them over its own priors under controlled distribution shift.
We build a synthetic tabular testbed in which covariate shift, spurious-feature reversal and concept shift (realised by giving features names whose everyday meaning agrees or disagrees with the rule) can each be switched on alone, with known ground truth.
On Qwen2.5-7B-Instruct with eight demonstrations, we compare seven hand-designed selection strategies against a zero-shot baseline.
We then train a set-transformer selector (\sata) on targets computed from the ground truth, and look inside the model with the logit lens, linear probes and activation steering.
Every confirmatory test was written down before its data were seen.

\textbf{Feature names, not demonstrations, decide most of what the model predicts.}
Flipping the names costs 0.20 \AUROC (95\% CI [0.13, 0.28], 22 of 24 tasks), more than any shift of the data.
No hand-designed strategy overcomes it.
\textbf{A trained selector is the only method with no weak environment on a stand-in learner, but it does no better on the instruction-tuned model}: none of the six pre-registered \sata contrasts favours the selector after Holm correction.
In an exploratory comparison, \sata does help the base model, whose priors pull less.
\textbf{Inside the model, the answer forms only in its last few layers, and that is where the names take over}: its zero-shot guess is decodable from layer 2, the demonstrated rule hardly at all, and a steering direction for ``follow the demonstrations'' moves its outputs without making them more correct.

Within the scale we could study (7B models, eight numeric demonstrations, synthetic tasks, a small compute budget), none of the selection methods we tried overcame the names' pull on the instruction-tuned model.
The testbed, the contrast families and the diagnosis give future work at larger scale a controlled place to test what can.

\cleardoublepage
\phantomsection
\chapter*{Acknowledgements}

\todo{Acknowledgements: supervisor, School of Computer Science, anyone who gave feedback, compute providers.}


% tables
\cleardoublepage
\setcounter{tocdepth}{1}
\tableofcontents

{\makeatletter
\renewcommand*\numberline[1]{\hb@xt@\@tempdima{#1 \hfil}\hspace*{1em}}
\makeatother
\listoffigures
\listoftables
\cleardoublepage
}

%%%%%%%%%%%%
% Chapters
\setcounter{page}{1}
\setcounter{chapter}{0}

% main page numbers:  1, 2, 3, ...
\renewcommand{\thepage}{\arabic{page}}
\setupParagraphs

\thesispart{Part I\quad Foundations}
\chapter{Introduction}
\label{ch:intro}

Large language models (LLMs) have shown a remarkable ability to learn classification tasks without being specifically trained for them, and hence without changing their internal weights. By placing a few labelled examples (\emph{demonstrations}) in the prompt, the model can be seen to answer the query prompt sufficiently well \citep{brown2020language, dong2024survey}. This phenomenon, known formally as \emph{in-context learning}, has been widely used by both users and professionals during training, testing and inference, largely due to its ease of use and setup. However, due to its fragile nature, it brings about challenges and varying results depending on which demonstrations are shown, their format and order, and their label balance \citep{zhao2021calibrate, lu2022fantastically, min2022rethinking}. Hence, a body of work arose to investigate the \emph{selection} of demonstrations: by similarity to the query \citep{liu2022what}, by diversity and coverage \citep{levy2023diverse,gupta2023coverage}, or with retrievers trained on the model's own feedback \citep{rubin2022learning,li2023unified,ye2023compositional}.


To the extent of the work that follows here, there are two things in practice that have been overlooked by most of the current work out there. 
The first is \textbf{distribution shift}: the queries a deployed model sees are rarely drawn from the same distribution as the demonstrations.
The second is \textbf{prior conflict}: a model brings its own beliefs about what a feature means, and what the demonstrations teach can disagree with them. Work by \citet{wei2023larger} found that following in-context evidence that contradicts semantic priors emerges with scale, and models of 7 to 8B parameters follow contradicting demonstrations only partially (not being able to override their priors and learn from what the demonstrations teach) \citep{krishnakumar2025semantic,kossen2024incontext}.


\newpage
\section{The question}

This thesis aims to ask a single question:

\begin{quote}
\textit{Can the choice of demonstrations help a small language model follow in-context evidence over its own priors, under controlled distribution shift?}
\end{quote}

The question is then broken into four research sub-questions, each answered on Qwen2.5-7B-Instruct \citep{qwen2024qwen25} with $k = 8$ demonstrations.

\begin{description}[leftmargin=2.4em, style=nextline]
  \item[\RQ1] Does in-context learning degrade under covariate shift, spurious-feature reversal and concept shift?
  \item[\RQ2] Do different selection strategies help under different kinds of shift?
  \item[\RQ3] Does a model of this size rely on the features that actually determine the label, and why does its prior about them win when it conflicts with the demonstrations?
  \item[\RQ4] Does a learned selector (\sata) beat the hand-designed strategies?
\end{description}


It is difficult to answer these questions on real data, as real data and the benchmarks built on them mix several kinds of shift and give no ground truth for which features matter (\autoref{sec:bg-shift}). Therefore, the work that follows builds a \textbf{synthetic tabular testbed} (\autoref{ch:method}) where every task has three features that determine its label, a spurious ``shortcut'' feature that agrees with the label 80--90\% of the time, and distractors. The shifts in focus here, covariate shift, reversal of the shortcut and concept shift, can each be switched on alone. Concept shift is realised by \emph{naming}: the same task is shown with neutral names ($f_0, \ldots, f_9$), with everyday names whose meaning agrees with the rule (\emph{aligned}), or with the antonyms of those names (\emph{flipped}), so that the model's semantic prior points the wrong way.

\newpage
\section{Approach and the story of this thesis}

The experiments here form a chain of sorts, and each chapter of Part~II aims to start from the limitation of the one preceding it.

\begin{enumerate}[leftmargin=2em]
  \item \textbf{Shift, names and the prior} (\autoref{ch:shift}; \RQ1--\RQ3).
        We measure how each shift hurts and whether seven hand-designed strategies help.
        The names dominate: no strategy overcomes flipped names, and the model's ``bigger numbers mean 1'' habit explains much of what remains.
        \emph{Limitation:} hand-designed heuristics see only the data, not what the model makes of it (information about the model's internal bias via its priors).
  \item \textbf{Learning to select} (\autoref{ch:sata}; \RQ4).
        We train a set transformer, \sata (Shift Aware Task Adapter), to choose demonstrations that teach the rule rather than the known shortcut.
        On a stand-in learner it is the only method with no weak environment, but on Qwen-Instruct it does no better than balanced random demonstrations. In an exploratory comparison, \sata does help the base model.
        \emph{Limitation:} behaviour alone cannot tell us why none of the demonstration sets we tried moves the instruction-tuned model.

  \item \textbf{Inside the model} (\autoref{ch:mechanism}; \RQ3).
        Using the logit lens, we saw that the answer forms in the last few layers and that the names decide it there; linear probes show that the demonstrated rule barely reaches the residual stream, and steering towards the demonstrations shifts every answer without making it more correct.
        \emph{Significant limitation:} our scale: one 7B model, eight numeric demonstrations, synthetic tasks.
\end{enumerate}

\section{Contributions}

\begin{enumerate}[leftmargin=2em]
  \item \textbf{A controlled testbed for in-context learning under shift and prior conflict}, with each shift isolated, ground-truth feature roles, a naming factor matched on measured prior strength, and a statistical protocol (within-cell \AUROC, a hierarchical bootstrap over tasks, Holm-corrected contrast families frozen before the data) (\autoref{ch:method}).
  \item \textbf{A measured account of how the 7B model here trades demonstrations against priors}: names outweigh every data shift, and the model's numeric prior explains most of the cost of reversing the shortcut (\autoref{ch:shift}).
  
  \item \sata, our demonstration selector, trained on ground-truth targets rather than on the model's scores, with an oracle check that at first caught a target teaching the shortcut backwards, and a pre-registered evaluation on the 24 evaluation tasks, which showed that \sata was not significantly better than the strategies on the instruction-tuned model at our small scale (\autoref{ch:sata}).

  \item \textbf{Evidence from inside the model} on where the names take over and how little of the demonstrated rule is visible there, and a negative result for single-direction steering (\autoref{ch:mechanism}).
\end{enumerate}

\section{Summary of findings}

The findings are stated in bold throughout, numbered by chapter, and always for the conditions we tested.
The main ones are these below.

\begin{itemize}[leftmargin=1.5em]
  \item \textbf{Flipping feature names hurts more than any shift of the data}: 0.20 \AUROC, against 0.12 for reversing the shortcut and 0.02 (not significant) for covariate shift (\autoref{ch:shift}).
  \item \textbf{Strategies that neutralise the shortcut help under its reversal more than under covariate shift} (interaction +0.11, Holm $p = 0.024$), but no strategy helps under flipped names (\autoref{ch:shift}).
  \item \textbf{A trained selector that wins on a stand-in learner does not win on the model}; none of the six pre-registered selector contrasts favours it (\autoref{ch:sata}).
  \item \textbf{The answer forms in the last six of 28 layers, and the names decide it there}; the model's prior is decodable from layer 2, while the demonstrated rule hardly is (\autoref{ch:mechanism}).
\end{itemize}
\chapter{Background}
The goal of this chapter is to build on the five strands of work the thesis draws on, as follows:

\begin{itemize}
    \item in-context learning and what it learns (\autoref{sec:bg-icl}),
    \item the contest between priors and demonstrations (\autoref{sec:bg-priors}),
    \item distribution shift and shortcuts (\autoref{sec:bg-shift}),
    \item demonstration selection (\autoref{sec:bg-selection}), and
    \item the tools for looking inside a model of this size, and what they reveal about how demonstrations are used (\autoref{sec:bg-mech}).
\end{itemize}

\section{In-context learning}
\label{sec:bg-icl}
Work by \citet{brown2020language} showed that a sufficiently large language model can perform a new task from a few demonstrations in its prompt, with no gradient updates; \citet{dong2024survey} survey the area since. However, given current progress in this area, what models seem to learn from demonstrations is being actively contested in the community. \citet{min2022rethinking} found that if the demonstration labels are replaced with random ones, the cost is often surprisingly small, suggesting that the demonstrations mainly convey the input distribution, the label space and the format. The works of \citet{pan2023what} and \citet{xie2022explanation} matter here because they
explain why a model may follow its priors when it recognises a task whose demonstrations disagree with them. \citet{pan2023what} separate \emph{task recognition}, in which the demonstrations only signal a task the model already knows from pre-training and the model answers from that knowledge (so even incorrect labels work), from \emph{task learning}, in which the model acquires a new input--label mapping from the demonstrations themselves.
They find that task recognition is present even in small models and does not improve with scale or more demonstrations, whereas task learning emerges only in larger models and improves as demonstrations are added. \citet{xie2022explanation} offer a complementary view: in-context learning as implicit Bayesian inference, in which the demonstrations update the model's prior over latent concepts, so a strong prior needs strong evidence to be overturned. A 7B model shown eight demonstrations is therefore expected to rely mostly on recognition, and to follow its priors whenever the demonstrations contradict them.

There have been several controlled studies that train or probe models on synthetic function classes. \citet{garg2022what} showed that transformers trained from scratch can learn linear functions and decision trees in context, and \citet{zhang2024trained} found that such in-context learners of linear functions tend to be brittle under covariate shift. In this same spirit, several tabular foundation models have been trained as well \citep{hollmann2023tabpfn, muller2022transformers, hollmann2025accurate}. Due to the autoregressive nature of LLMs, we follow a similar technique to \citet{brown2020language}, in which classification is read from the probabilities of label tokens, and we serialise tabular rows following \citet{dinh2022lift, hegselmann2023tabllm}; however, this method of reading the probability of tokens is affected by surface-form competition \citep{holtzman2021surface} and by bias towards particular labels. We aim to correct this via contextual calibration \citep{zhao2021calibrate}, and we avoid public datasets because of the open possibility that they were memorised during the training of models like Qwen \citep{bordt2024elephants}.

\section{Priors against demonstrations}
\label{sec:bg-priors}

A model's \emph{prior} is defined as what it would predict with no demonstrations. When feature names carry meaning, the prior is semantic: a feature called ``credit score'' suggests a direction before any example is seen. \citet{dinh2022lift} and \citet{choi2022lmpriors} have shown this effect where feature names change what the language model does with tabular data. 

\citet{wei2023larger} tested what happens when the demonstrations contradict the prior, by flipping labels or replacing them with semantically unrelated ones: only large models follow the flipped labels, and smaller models stay with their priors. At the scale of this study, we found this to be partly true. Further studies by \citet{krishnakumar2025semantic} have shown that models like Llama-3.1-8B-Instruct agree with fully flipped sentiment labels about half the time, and while the statistical notion of prior--data conflict goes back to \citet{evans2006checking}, related work on knowledge conflicts \citep{longpre2021entity, xie2024adaptive} studies when a model trusts its context over its parametric memory. This thesis treats the question as a contest between prior and demonstrations whose outcome can be shifted, and asks whether the demonstrations themselves can be chosen to shift these outcomes.

\section{Distribution shift and shortcuts}
\label{sec:bg-shift}

\paragraph{Kinds of shift.}
Covariate shift changes the input distribution $P(X)$ while keeping $P(Y \mid X)$; label (prior-probability) shift changes $P(Y)$; concept shift changes $P(Y \mid X)$ itself \citep{shimodaira2000improving,quinonero2009dataset,morenotorres2012unifying}.
On real tabular benchmarks, most of the large performance drops are dominated by concept shift \citep{liu2023need}.
Benchmarks such as TableShift \citep{gardner2023tableshift} collect natural shifts, but each mixes several kinds, so the effect of one kind cannot be isolated.

\paragraph{Spurious features and shortcuts.}
A spurious feature is a feature with strong predictive power that predicts the label in the training environment for reasons that do not hold elsewhere, particularly in real-world situations, and it has been shown that learners adopt such a feature as a shortcut for prediction \citep{geirhos2020shortcut, nagarajan2021understanding}. In-context learning inherits the problem: when demonstrations are consistent with several features, models show systematic preferences among them \citep{si2023measuring} and exploit shortcuts in the prompt \citep{tang2023large}, and in-context learners trained on tasks with spurious features adopt them \citep{harutyunyan2026incontext,chen2026incontext}. Restricting a predictor to only causal features is not automatically better either: on tabular shifts, \citet{nastl2024causal} found that predictors using all the features generalised as well or, at times, better.

\section{Demonstration selection}
\label{sec:bg-selection}
\paragraph{Heuristic selection.}
The strategies we build for demonstration selection follow general rules of thumb that the research community has worked on over the years. The simplest of these is retrieving the training examples nearest to the query \citep{liu2022what}, while others cover the input space with diversified examples, show demonstrations that make the shortcut uninformative, show the model examples that correct its mistakes, or control label balance.

\newpage
\paragraph{Learned selection.}
Several works train a selector based on a model's own feedback. EPR \citep{rubin2022learning} and UDR \citep{li2023unified} train a separate retriever to choose demonstrations: each candidate is scored by how much it raises the language model's probability of the correct answer, and the retriever learns to rank the most helpful candidates first, while others like CEIL \citep{ye2023compositional} learn to choose whole sets (grouping examples), and \citet{zhang2022active} learn a selection policy by reinforcement learning. Techniques ranging from regressing a model's output on which training examples were present \citep{ilyas2022datamodels} to learned judges built from a model's internal signals \citep{kato2025affinity, wang2026disp} have also been proposed; \citet{kato2025affinity} show that scores computed from induction-head representations track in-context accuracy, but leave selection with them to future work.

\paragraph{Set-based architectures.}
The selector built in this thesis reads the whole pool of candidates at once.
It draws on In-Context Risk Minimization \citep{gupta2024context}, in which a transformer adapts to an environment by attending over other inputs from it, on the style of TabPFN, where the tokenisation is done at the row level of the table \citep{hollmann2023tabpfn}, and on attention over sets \citep{lee2019set}.

\section{Looking inside the model}
\label{sec:bg-mech}
The running representation that every layer reads from and writes to can be observed in the residual stream of the transformer \citep{vaswani2017attention, elhage2021mathematical}. The \emph{logit lens} reads an intermediate layer's prediction by applying the final normalisation and unembedding to it \citep{nostalgebraist2020logit}; its readings can be biased at early layers \citep{belrose2023eliciting}.
\citet{halawi2024overthinking} used it to find a \emph{critical layer} after which models given false demonstrations switch to the demonstrated labels. Linear probing tests whether information is linearly decodable from a layer \citep{alain2017understanding}, but shows only that the information is present, not that it is used \citep{belinkov2022probing}. To change behaviour by adding a direction to the residual stream, one can take the difference of mean activations, as in activation addition \citep{turner2023activation} and contrastive activation addition \citep{rimsky2024steering}. Other work has found a single direction that controls whether a model follows its context or its priors \citep{minder2025controllable}.

\section{The gap}
\label{sec:bg-gap}

Three gaps meet in this thesis.
First, selection methods are almost always evaluated in distribution, with no control over what kind of shift the queries undergo.
Second, the conflict between semantic priors and demonstrations has been studied by flipping \emph{labels}, not by changing what the \emph{features} mean, and not under shift. Third, the learned selectors above are trained on the model's scores of whole prompts; none, as far as our search reached, uses the model's measured prior on each candidate, and none is tested on whether it helps a small model overcome that prior. The nearest work, \citet{kato2025affinity}, is correlational and builds no selector. This thesis builds the testbed that the first two gaps require, and uses it to test the third.

\chapter{Method: a controlled testbed}
\label{ch:method}

Every experiment in Part~II runs on the same testbed: synthetic tasks with known structure, three environments that each change one thing, a naming factor that sets the model's prior against the rule, a fixed prompt and scoring procedure, and a shared statistical protocol.
This chapter describes each in turn. Every symbol is listed in \autoref{app:notation}, and the full definitions, with proofs of the properties stated here, are in \autoref{app:design}.

\section{Why synthetic tasks}
\label{sec:why-synthetic}

The project began on real tabular benchmarks.
TableShift \citep{gardner2023tableshift} names covariate, label and concept shift, but every one of its shifts mixes them, so a drop in accuracy cannot be attributed to one kind.
Real data also give no ground truth for which features determine the label, which \RQ3 needs, and public datasets may have been memorised \citep{bordt2024elephants}.
Synthetic tasks solve all three problems, in the tradition of \citet{garg2022what} and of the synthetic priors behind tabular foundation models \citep{muller2022transformers,hollmann2023tabpfn}.

\section{Tasks}
\label{sec:tasks}

Each task has ten features.
Three are \emph{load-bearing}: the label depends on them and on nothing else.
Five are \emph{distractors}.
One, $f_8$ in the generator's numbering, is a \emph{spurious shortcut}, and one, $f_9$, is pure noise.
The ten features are shown in a random order recorded per task, so no feature's role can be read from its position.

\paragraph{Rules.}
Half of the tasks are \textbf{linear}: $P(y = 1 \mid x) = \sigma(c^\top x_C)$ over the load-bearing features $C$, with coefficients $c_j = u_j r_j$, $u_j \sim \mathrm{U}(1,2)$ and a random sign $r_j$, so every load-bearing feature has a definite direction.
The other half are \textbf{trees}: a signed majority vote over the three features' signs, which makes them monotone in each feature as well.

Labels are then flipped with probability $\eta = 0.02$. A row whose label was not flipped is \emph{clean}, and the rule's label is the \emph{clean label} $y_{\mathrm{clean}}$. This \emph{label noise} is separate from $f_9$, the noise feature, which is an input that has no bearing on the label.

\paragraph{The shortcut.}
The spurious feature is generated from the clean label,
\[
  x_8 = \sigma_t\big((2y_{\mathrm{clean}} - 1)\,\Phi^{-1}(s) + \varepsilon\big), \qquad \varepsilon \sim \mathcal{N}(0,1),
\]
so its sign agrees with the label with probability $s \sim \mathrm{U}(0.80, 0.90)$, the range of colour--label agreement in Colored MNIST's training environments \citep{arjovsky2019invariant}.
On its own it predicts the label better than any single rule feature.
The spurious direction $\sigma_t \in \{+1, -1\}$ is drawn per task and balanced, for a reason that only appeared once the model was run: Qwen's zero-shot answers rise with every feature value, so with $\sigma_t$ fixed at $+1$ that habit backed the shortcut in every task (\autoref{sec:shift-p2}).

\paragraph{Suites.}
The evaluation suite has 24 tasks: 12 linear and 12 tree, crossed with two domains (loan approval and medical risk) and balanced on $\sigma_t$.
A separate 12-task pilot suite was used to choose the number of demonstrations, so that choice never touched the evaluation tasks.
Each task has a 256-row demonstration pool and, per environment, a label-balanced set of 20 queries.
All features are standardised with the pool's statistics.
The selector of \autoref{ch:sata} is trained and validated on further tasks from the same generator, drawn with different random seeds so that it never sees an evaluation task.

\section{Environments}
\label{sec:environments}

Queries come from three environments; the demonstration pool is always in distribution.

\begin{description}[leftmargin=1.5em, style=nextline]
  \item[In distribution (ID)] Queries are drawn exactly like the pool.
  \item[Covariate shift] Two load-bearing features are moved by 1--2 standard deviations in opposite label directions, with equal pushes on the label, and one distractor is moved as well.
        $P(y \mid x)$ is unchanged, and so is the label rate, by construction.
        An earlier design filtered random shifts on the label rate; it turned out to accept almost only shifts of distractors, which cannot affect any label (21 of 24 tasks).
  \item[Spurious reversal] The shortcut's sign is reversed: it now agrees with the label in only $1-s$ of rows.
\end{description}

\textbf{Concept shift} is realised by naming (next section), because changing $P(y \mid x)$ itself while keeping the task recognisable is exactly what flipped names do from the model's point of view.

\section{Naming: setting the prior against the rule}
\label{sec:naming}

Each task is shown under three namings.
\textbf{Abstract} names are $f_0, \ldots, f_9$.
\textbf{Aligned} names give each load-bearing feature a domain name whose everyday meaning agrees with its direction in the rule: in the loan domain, a feature that raises approval might be called ``credit score''.
\textbf{Flipped} names give the same feature the antonym, ``missed payments''.
Aligned and flipped prompts differ only in the three load-bearing names; non-causal features get plausible names with weak priors (``form version'', ``office floor number''), so names do not reveal which features matter.

Antonym pairs were screened on the model itself.
On 2,000 zero-shot profiles, a numeric prior dominated: every abstract feature had a positive slope, and names that should lower $P(y = 1)$ mostly just raised it less.
Only 1 of 16 candidate pairs met the planned symmetry rule.
The final names are the five pairs per domain with the largest \emph{measured} contrast.
On the pilot tasks, zero-shot \AUROC is 0.594 with aligned names and 0.430 with flipped names ($\Delta = 0.164$, 95\% CI [0.053, 0.287], 10 of 12 tasks), so the names do move the prior; the effect is large for loan (0.28) and weak for medical (0.05).

\section{Prompts, scoring and the model}
\label{sec:prompts}

Rows are serialised as \texttt{name: value} pairs separated by semicolons, followed by \texttt{->} and the label.
The system message states that each example lists ten standardised measurements and a label, that the label follows an unknown rule over some of them, and that the model should answer with one character, 0 or 1.
An abstract-name prompt reads, in part:

\begin{quote}\small\ttfamily
f0: 0.99; f1: 1.42; f2: 0.72; f3: -0.23; f4: -0.52; f5: 1.09; f6: 0.35; f7: -0.74; f8: -0.99; f9: -1.63 -> 0\\
\ldots\\
f0: 1.25; f1: 0.19; f2: 0.10; f3: -0.04; f4: -0.83; f5: 0.20; f6: 0.01; f7: 0.36; f8: -0.31; f9: 0.65 ->
\end{quote}

The model is \textbf{Qwen2.5-7B-Instruct} \citep{qwen2024qwen25} in bf16, run through PyTorch and Hugging Face \texttt{transformers} \citep{paszke2019pytorch,wolf2020transformers}.
Its score for a query is the margin $m = \log p(\texttt{1}) - \log p(\texttt{0})$ of the two label tokens, computed in fp32 from the final hidden state; reading it from the bf16 output layer rounded margins onto a coarse grid.
The demonstration prefix is computed once and cached, and each query is scored against the cached keys and values \citep{pope2023efficiently,gim2024prompt}; this was checked to change mean per-task \AUROC by at most 0.005.
Qwen2.5-7B (base) answers in format and is used for exploratory comparisons. Inference ran first on an Apple M3 Max and then on a rented NVIDIA H200 with eight copies of the model sharing the GPU (the main grid, about 255,000 forward passes, took 42 minutes); the whole project used a few tens of US dollars of GPU time, which bounds what could be tested (\autoref{sec:limitations}).

Llama-3.1-8B-Instruct \citep{grattafiori2024llama3} was dropped after the pilot, because with this prompt it explains its reasoning instead of answering (only about 2\% of its probability fell on the label tokens).

\paragraph{The number of demonstrations.}
On the pilot suite, $k = 8$ gold demonstrations beat the same rows with shuffled labels by $\Delta\AUROC = 0.098$ (95\% CI [0.021, 0.188], 9 of 12 tasks), so the model does use the labels; $k = 32$ did not pass.
All experiments use $k = 8$.

\paragraph{Measured priors.}
The model's prior on a row is its zero-shot margin on that row, under the unit's naming, centred on the pool median.
A ridge surrogate of this prior from feature names \citep{hoerl1970ridge} failed its planned fit criterion (held-out $R^2$ 0.20--0.51 against a bar of 0.5), so selection uses the measured margins directly.

\section{Selection strategies}
\label{sec:strategies}

\autoref{tab:strategies} lists the strategies.
Every strategy except \strat{random} takes $k/2$ demonstrations of each label, and balanced strategies are compared with \strat{label\_diversity}, a label-balanced random draw, so that contrasts compare selection mechanisms at fixed label counts.

\begin{table}[t]
\centering\small
\caption{Selection strategies. Ground-truth strategies (marked GT) see feature roles that a deployed selector would not; this is disclosed wherever they are compared with \sata.}
\label{tab:strategies}
\begin{tabularx}{\textwidth}{lX}
\toprule
\textbf{Strategy} & \textbf{What it selects} \\
\midrule
\strat{zero\_shot} & No demonstrations. \\
\strat{random} & Uniform draw from the pool. \\
\strat{label\_diversity} & Uniform draw, $k/2$ rows of each label: the comparator. \\
\strat{feature\_range} & Covers the quantile bins of the three features with the highest mutual information with the label. \\
\strat{rule\_diversity} (GT) & Covers the ground-truth regimes of the rule, one row each. \\
\strat{counter\_spurious} (GT) & Matched pairs on $|x_8|$ in which the shortcut agrees and disagrees equally often, so it carries no information. \\
\strat{similarity} & The $k/2$ nearest rows of each label to the query in feature space \citep{liu2022what}. \\
\strat{counter\_prior} & Clean rows whose label contradicts the model's measured prior, when the pool shows prior--data conflict; related to misconfidence-based selection \citep{xu2024misconfidence}. \\
\sata & Learned selector (\autoref{ch:sata}). \\
\bottomrule
\end{tabularx}
\end{table}

\section{Evaluation and statistics}
\label{sec:stats}

\paragraph{The metric.}
The primary statistic is \AUROC \citep{hanley1982meaning,fawcett2006introduction} of the margin, computed within each cell of (task, strategy, seed, naming, environment) over its 20 label-balanced queries.
It was chosen over calibrated balanced accuracy, after the first tasks had run and before any contrast was computed: contextual calibration left every query of a cell on one side of the threshold in 36\% of plain-random cells, which pins balanced accuracy at 0.5 and forces shift gaps to zero.
Within-cell \AUROC has a blind spot of its own (\autoref{sec:limitations}): a shift that moves every score together cannot be seen.

\paragraph{Inference.}
For most strategies the demonstration set depends only on (task, strategy, seed), so the same sets are reused across namings and environments, and contrasts are paired; \strat{similarity} and \sata choose a set for each query, and \strat{counter\_prior}'s set depends on the naming, because it uses the prior measured under that naming.
Uncertainty comes from a \textbf{hierarchical bootstrap} \citep{efron1993introduction,davison1997bootstrap}: each of 2,000 draws resamples the 24 tasks and then, within each task, the positive and negative queries separately.
Contrasts are one-sided, and each family is \textbf{Holm-corrected} \citep{holm1979simple}.
Every confirmatory family was \textbf{written into the code before the data it tests were seen} \citep{nosek2018preregistration}: ten contrasts for \RQ1--\RQ3, six for \sata.
Everything outside those families is labelled exploratory.

\paragraph{Figures.}
Following \citet{healy2019data}, estimates are positions on a common scale with 95\% intervals, categories are ordered by value, and colour carries one meaning per figure, in colour-blind-safe palettes; every figure is drawn from the result files by one script.

\paragraph{Gates.}
Each expensive stage was preceded by a check that the instrument worked (cache exactness, answer format, learnability, naming strength, and an oracle check of every selector's targets), because an earlier, ungated version of this project produced results that turned out to be artefacts of code defects.


\thesispart{Part II\quad Experiments}
\chapter{Shift, names and the prior}
\label{ch:shift}

\emph{Research questions: \RQ1, \RQ2 and the behavioural half of \RQ3.}

\section{Motivation}
Before asking whether a selector can help, we needed to know what was going wrong: how each shift affected the model tested here, Qwen2.5-7B-Instruct; whether a hand-designed strategy helps under the shift it was designed for; and which features the model relies on.

\section{What the pilot revealed: a numeric prior}
\label{sec:shift-p2}
In the first pilot run the shortcut's direction remained the same across all tasks: a larger value of $x_8$ always meant the label was 1. For the pilot ID queries under abstract names, we calculated the raw sum of the features $f_0$ through $f_9$ and compared it with the margin $m = \log p(\texttt{1}) - \log p(\texttt{0})$ of Qwen's prediction for the same query using Spearman's correlation, which was 0.70. This meant that Qwen's internal approach carried a habit, a prior, similar to the summing approach. Scoring those same ID queries by that sum gives \AUROC 0.628, and 0.511 once $x_8$ is left out; Qwen's own zero-shot \AUROC was 0.633. Its ``zero-shot skill'' was the shortcut, read through a habit that bigger numbers mean 1.

\finding{Qwen2.5-7B-Instruct carries a numeric prior, ``bigger values mean 1'', that applies to every feature, including the shortcut.}

\newpage
It was interesting to see that this model followed its priors, essentially its gut, across zero-shot tasks, with an \AUROC of 0.642 where the habit backs the shortcut ($\sigma_t = +1$) and 0.357 where it opposes it ($\sigma_t = -1$), despite the shortcut's direction being made random per task and balanced. \citet{wei2023larger} show that the ability to do ICL well and override priors appears with scale; this is a limitation we try to address with demonstrations and, later, with a selector trained to choose demonstrations that teach the rule.

\begin{figure}[H]
    \centering
    \includegraphics[width=\textwidth]{\figpath/fig_rq1_rq3_contrasts.png}
    \caption{Five of the ten confirmatory contrasts: the three for \RQ1 and the two for \RQ3 (the five \RQ2 contrasts are in \autoref{tab:contrasts}). Each row is an estimate with its 95\% interval (hierarchical bootstrap over 24 tasks); filled blue marks the two that survive Holm correction (adjusted $p$ at the right, corrected over all ten), hollow grey the rest. The three panels use different units, so they do not share a scale. Flipping the names (concept shift) is the largest effect.}
    \label{fig:p4-contrasts}
\end{figure}

\section{Setup}

The experiments use the 24 evaluation tasks of \autoref{sec:tasks}: 12 linear and 12 tree tasks, generated synthetically across the loan and medical domains. Each has a 256-row in-distribution pool from which demonstrations are drawn, and 20 label-balanced queries in each of three environments (ID, covariate shift and spurious reversal; \autoref{sec:environments}).
The \emph{main grid} is the full set of runs over every combination of task, naming (abstract, aligned, flipped) and condition (zero-shot and the seven hand-designed strategies of \autoref{tab:strategies}), with three demonstration seeds per strategy (one for \strat{similarity}, whose nearest-neighbour sets do not depend on the seed).
In each run the strategy picks eight demonstrations from the pool, and that same set is paired with each of the task's 60 queries, giving one margin per query and one \AUROC per environment from its 20 queries.
Except for \strat{similarity}, which picks a set for each query, and \strat{counter\_prior}, whose set depends on the naming, a demonstration set depends only on the task, strategy and seed, so the same rows are shown under every naming and environment, and differences between them are never differences in the demonstrations.
This gives 86,400 scored prompts on Qwen2.5-7B-Instruct, and a reduced grid on Qwen2.5-7B (base) with abstract and flipped names only.

\section{\RQ1: How much does each shift hurt?}

\finding{Flipping the feature names costs more than any shift of the data: 0.20 difference in \AUROC for random demonstrations on ID queries (95\% CI [0.13, 0.28], 22 of 24 tasks).}

Based on the results in \autoref{fig:p4-contrasts}, we can see that the most significant cost was for flipping names. The basis of the metric, the \emph{difference in AUROC}, is a paired comparison made within each task. For every one of the 24 tasks we take Qwen's margins on 20 label-balanced queries in two conditions that differ in exactly one respect, compute the \AUROC in each (the chance that a randomly chosen label-1 query receives a higher margin than a randomly chosen label-0 one), subtract, and average the 24 differences; the interval comes from the hierarchical bootstrap of \autoref{sec:stats}. A positive value is the \AUROC lost to the change. In the first two rows the change is the environment, ID queries against covariate or reversal queries, both under abstract names so that no name meaning is involved; in the third it is the naming, the same ID queries under aligned and under flipped names.

The third row is how we realise concept shift, and it cannot be built like the other two: if the queries followed a different rule from the demonstrations, the demonstrations would simply be wrong and no choice of them could help. Flipped names keep the data, the rule and the demonstrations fixed and move the conflict to the model's side instead, because the rule it brings from pre-training (a higher credit score means approval) now points the opposite way to the task's. From the model's point of view the mapping from features to label has changed, which is what concept shift means for a learner whose prior is semantic, even though $P(y \mid x)$ in the data never does.

Read this way, the ordering of the three rows is the main result of this section. In the flipped condition the prompt carries exactly the evidence it carries in the aligned one, the same numbers, the same eight correctly labelled demonstrations and the same queries, and only three feature names differ, yet Qwen loses 0.20 \AUROC, and in 22 of the 24 tasks. Eight demonstrations that show the correct mapping are not enough to overrule what the names mean to it. Nor is this a weakness of random demonstrations: on ID queries every strategy falls to 0.44--0.54 under flipped names, against 0.66--0.71 under aligned ones (\autoref{fig:strategies-by-naming}).

This is the behaviour \citet{wei2023larger} describe for models below the scale at which in-context evidence overrides semantic priors, and a 7B model sits below it. Qwen does not ignore its demonstrations, though. With \strat{label\_diversity}'s sets, correct labels beat shuffled ones on ID queries under every naming, by 0.06 to 0.11 \AUROC, and under flipped names they lift Qwen from 0.41 only to 0.51. The demonstrations cancel the names' pull rather than overturn it, in line with \citet{kossen2024incontext}, who find that in-context labels change a model's predictions without fully overcoming the preferences it acquired in pre-training.

Reversing the shortcut costs about half as much (+0.12), as seen in \autoref{fig:p4-contrasts}. The interval excludes zero and the direction holds in 16 of 24 tasks, but the effect misses significance after Holm correction ($p = 0.063$), and \autoref{sec:shift-prior-shortcut} shows why it is so noisy: with random demonstrations under abstract names the cost is 0.21 in the tasks where the ``bigger means 1'' habit backs the shortcut and 0.02 where it opposes it, so the pooled estimate averages a large effect with almost none. Much of what looks like a cost of reversal is the numeric prior of \autoref{sec:shift-p2} reading the shortcut in the query, rather than something the demonstrations taught.

Covariate shift costs nothing we can measure (+0.02, 15 of 24 tasks). Its interval runs from $-0.04$ to 0.09, so what the data rule out is a large cost, not a moderate one, and the shift itself is not small: it moves two of the three rule features by one to two standard deviations. However, moving where the inputs fall disturbed Qwen far less than changing what the inputs are called.

For this 7B model, then, the threat to in-context learning is less a shift in the data than a conflict between the task and what the model already believes. If that is so, better demonstrations can help only as far as the model is willing to read them, which is what \RQ2 tests next. The two \RQ3 rows of \autoref{fig:p4-contrasts} are taken up in \autoref{sec:shift-rq3}.

\section{\RQ2: does each strategy help under its target shift? (strategy - label diversity)}
\label{sec:shift-rq2}

\finding{The strategies' value depends on the kind of shift, and they help at most marginally (not significantly) in situations where the semantic relationship in the features is reversed or lost.}

Comparing random demonstrations and each strategy under the two shift types, covariate shift and shortcut reversal, against the baseline of ID test rows within a task, we saw that the strategies work best when the feature names are aligned with a clear semantic meaning, and worst when they are flipped.

Across the ID environment we can see that the counter-prior, rule-diversity and similarity strategies performed slightly above chance under all naming schemes. Under the shifts, the similarity strategy does well under covariate shift (the best of all under flipped names) and worst, or close to worst, under spurious reversal, meaning that this strategy largely emphasised the predictive power of the shortcut. Conversely, the counter-spurious strategy performs best under spurious reversal, where the shortcut points the wrong way.

Overall, what we take away is that, across strategy designs, attempting to override a model's priors gave the highest mean \AUROC when the features carry semantic meaning, although its lead over \strat{label\_diversity} is small and not significant. It is possible that combining this strategy with the counter-spurious strategy could help to deal with models' reliance on shortcut features.

It is worth mentioning that the rule-diversity strategy was observed to perform nearly as well as the attempt to reverse the priors under every naming, making it a possible substitute, although it reads the ground-truth regimes of the rule, which a deployed selector would not know.

\begin{figure}[H]
    \centering
    \includegraphics[width=\textwidth]{\figpath/fig_strategies_by_naming.pdf}
    \caption{Every strategy by naming (panels) and environment (dot colour), as mean AUROC over 24 tasks. The grey line joins one strategy's three environments, so its length shows how much the shifts move that strategy. Rows are sorted by their mean under aligned names. The panels differ far more than the rows within a panel: the names matter more than the choice of demonstrations.}
\label{fig:strategies-by-naming}
\end{figure}

\begin{table}[H]
\centering\small
\caption{The five planned \RQ2 contrasts: each strategy against \strat{label\_diversity} under the shift it targets.}
\label{tab:rq2-rows}
\begin{tabularx}{\textwidth}{@{}l l X@{}}
\toprule
\textbf{Row} & \textbf{Strategy tested} & \textbf{Its target shift} \\
\midrule
counter\_spurious, under reversal & \strat{counter\_spurious} & The shortcut reversing \\
similarity, under covariate & \strat{similarity} & Covariate shift \\
feature\_range, under covariate & \strat{feature\_range} & Covariate shift \\
counter\_prior, under flipped names & \strat{counter\_prior} & Names pointing the wrong way \\
counter\_spurious's gain: reversal $-$ covariate & \strat{counter\_spurious} & Does its help depend on the shift? \\
\bottomrule
\end{tabularx}
\end{table}

\section{\RQ3: which features does the model rely on?}
\label{sec:shift-rq3}
To find out whether Qwen relies on the true features, and whether it uses them in the direction the demonstrations show or the direction the names suggest, we ran tests on the 20 ID queries of each of the 24 tasks under abstract, aligned and flipped feature names. The \emph{label diversity} strategy was used to construct the demonstrations here because it provides the balanced targets needed for prediction. To see which of the 10 features Qwen uses ($\pi_{behav}$), we use a removal-based measure of importance \citep{breiman2001random,covert2021explaining}: each feature is replaced by hot-deck imputation \citep{andridge2010review} with a value from the same pool of rows the demonstrations are drawn from, which keeps the altered query in distribution \citep{zhu2026faithfulness}, and we observe the change in Qwen's resulting answer through the margin, computing the \emph{drop} as the correct-label margin (original) $-$ the correct-label margin (swapped). For the bias correction, the code also subtracts a ``content-free'' baseline, obtained by passing in a query whose feature values are \emph{N/A} \citep{zhao2021calibrate}.

From the construction of the tasks we know $\pi_{true}$, the ground-truth features the generator constructed, also known as the rule features; synthetic tasks give this kind of ground truth for judging feature importance \citep{liu2021synthetic}.

\finding{With abstract names, Qwen follows the demonstrations slightly more often than chance but barely learns the rule's direction from the examples. With flipped names, which conflict with the rule, and despite the demonstrations showing the correct direction, Qwen's agreement falls to 42\%, meaning it goes with the names more often than with the evidence in front of it. The names do not just add noise; they actively pull Qwen's use of each feature the wrong way.}

The \emph{demo-following index} was used to test whether the model indeed followed the demonstrations it was given. As a sensitivity test, this was done by nudging each load-bearing feature by $\pm 0.5$ and measuring how often the margin moved in the rule's direction; we saw that the model does follow the demonstrations, as shown in \autoref{fig:p4-contrasts}. We compare with abstract rather than aligned names because under aligned names following the demonstrations and following the names move the answer the same way, so the two cannot be told apart; however, it is worth noting that Qwen's prior bias, ``bigger means 1'', is still active even under abstract names.

\finding{With abstract names, Qwen relies on the shortcut, not the rule, but with semantically aligned, meaningful features we can see the model go in the right direction, relying on the rule features.}

In assessing the correlation of $\pi_{self}$ and $\pi_{behav}$, a comparison of self-reported and behavioural rankings in the manner of \citet{li2025stads}, all 216 rankings were parsed. With names, Qwen ranks the named rule features in its top three (0.64--0.69 of rankings) even under flipped names, and the rank correlation between what it reports and what it relies on is $-0.02$ (abstract), 0.12 (aligned) and $-0.02$ (flipped), in line with work on unfaithful self-explanations \citep{huang2023can,turpin2023language,lanham2023measuring}.

\begin{figure}[H]
\centering
\includegraphics[width=0.78\textwidth]{\figpath/fig_pi_behav.pdf}
\caption{Which features Qwen relies on ($\pi_{\mathrm{behav}}$, the drop in correct-label margin when a feature is replaced by a value from another row), by naming. Under abstract names the shortcut carries the reliance; under aligned names the rule features do; under flipped names the rule features are used in the names' direction (negative).}
\label{fig:pi-behav}
\end{figure}

\section{The prior reads the shortcut in the query}
\label{sec:shift-prior-shortcut}

Splitting the cost of spurious reversal by the shortcut's direction explains much of it (\autoref{fig:reversal-direction}).
Where the ``bigger means 1'' habit agrees with the shortcut ($\sigma_t = +1$), reversal costs every strategy 0.04--0.21 \AUROC.
Where it opposes the shortcut ($\sigma_t = -1$), most gaps are near zero or negative.
Zero-shot, with no demonstrations at all, splits the same way: $+0.07$ to $+0.13$ against $-0.09$ to $-0.15$.

\begin{figure}[t]
\centering
\includegraphics[width=\textwidth]{\figpath/fig_reversal_direction.pdf}
\caption{The cost of spurious reversal (ID minus reversal \AUROC), split by whether Qwen's ``bigger means 1'' prior agrees with the shortcut ($\sigma_t = +1$, vermilion) or opposes it ($\sigma_t = -1$, blue). The cost appears where the prior agrees with the shortcut, with or without demonstrations (zero-shot row).}
\label{fig:reversal-direction}
\end{figure}

\paragraph{Demonstrations without the shortcut.}
If demonstrations are meant to teach the causal rule, perhaps they should not show the shortcut at all.
As an exploratory experiment we removed the shortcut (and, in a later stage, the noise feature) from every demonstration while keeping queries complete.
Removing the shortcut cut ID \AUROC by 0.048 and covariate \AUROC by 0.040, and raised spurious-reversal \AUROC by 0.049.
The ID--reversal gap shrank in all 18 (strategy, naming) cells, significantly in 14.
Removing the shortcut worked as well as neutralising it with \strat{counter\_spurious}, and removing the noise feature added nothing.
The gap that remains follows $\sigma_t$ ($+0.04$ to $+0.11$ against $-0.06$ to $-0.15$): the prior still reads the shortcut in the query.

\finding{Much of what looks like shortcut learning from demonstrations is the model's numeric prior reading the shortcut in the query; hiding the shortcut from the demonstrations removes the demonstration part and leaves the prior part.}

Hence, it is possible that the model maps the query task to a similar task it has seen before and prefers to copy the method for that task.

\paragraph{The base model.}
Qwen2.5-7B (base) loses less to flipped names than the instruction-tuned model (random demonstrations on ID queries: 0.625 to 0.560, against 0.603 to 0.484), which matters in \autoref{ch:sata}, where it is the one model a learned selector helps, in an exploratory comparison (\autoref{app:base}).

\newpage
\section{What this chapter does not show, and the next question}

Every strategy tested reads only the data, never what the model makes of it (\strat{counter\_prior} is a single-rule heuristic, inactive in nine tasks), and only one comparison, the interaction for \strat{counter\_spurious}, shows that matching strategy to shift pays; with 24 tasks, effects below about 0.05 \AUROC are hard to detect.
The natural next step is to \emph{learn} a selector that composes these ideas, choosing per query the rows that teach the rule rather than the shortcut: \autoref{ch:sata}.

\chapter{Learning to select}
\label{ch:sata}

\emph{Research question: \RQ4 (does a learned selector beat the hand-designed strategies?).}

\section{Motivation}

\autoref{ch:shift} left an opening. We saw that neutralising the shortcut fared better under its reversal than under covariate shift. But the strategies that did so were hand-built rules, each of which targeted only one shift at a time. This motivated us to explore the idea of a learned selector that combines these ideas, choosing per query the rows that teach the rule rather than the shortcut, and that can be trained on synthetic tasks where the best choices are known. In this chapter, we build \sata (\autoref{sec:sata}).


\section{The selector}
\label{sec:sata-arch}
\sata scores every row of the 256-row demonstration pool, given the whole pool and the query. We set each candidate to be a single token made up of its feature vector plus a label embedding plus a ``demonstration'' type embedding; the query is a token with its feature vector and a ``query'' type embedding. \sata is an encoder-only transformer ($d = 128$, 4 heads, 4 layers, about 0.8M parameters) that attends over all tokens, and a linear head gives one score per candidate, normalised by a softmax. There are no positional encodings, so the encoder remains permutation-equivariant over candidates \citep{lee2019set}; the row encoding follows TabPFN's style of approach \citep{hollmann2023tabpfn}, and the idea of adapting by attending over other inputs from the environment comes from \citet{gupta2024context}.

\sata is not trained on the language model's scores; its aim is to learn targets computed from the ground truth, so its training never touches the model it is evaluated on. This is the distinguishing difference in comparison to EPR, UDR or CEIL \citep{rubin2022learning,li2023unified,ye2023compositional}.

\paragraph{Getting the targets right.}
\label{sec:sata-targets}
\sata learns by example, and for each pool it needs a ``right answer''. This is essentially a score for every row saying how strongly that row should be picked. The answer key is the target here, and because we built the tasks synthetically, we know the ground truth, which makes our target known as well, computed outside Qwen. However, due to computation limits, we needed a way to test these answer keys before training \sata on them. The idea was an oracle check.

Each oracle picks 8 demonstrations, and we tested 3 oracles. The check did not use the 24 evaluation tasks. It ran on \sata's own 100 validation tasks, drawn from the same generator but kept separate, so the evaluation tasks stayed untouched until the final test on Qwen. Our first approach was almost greedy: it rewarded rows in the query's region and rows that break the shortcut, and penalised the ones that follow the shortcut and sit in a different region. The issue with this approach was that it favoured label-noise rows instead. In 2\% of rows the label is flipped by noise, but the shortcut was generated from the true label, so in those rows it disagrees with the recorded label. They look like shortcut-breakers, but their label is actually wrong, so the targets were teaching wrong examples. Upon discovering this, we excluded the label-noise rows, only to find that the targets then overshot: they picked so many shortcut-breaking rows that the shortcut agreed with the label in only 25\% of the chosen rows. That does not remove the shortcut; it just teaches it backwards. The solution we settled on was to copy the \strat{counter\_spurious} strategy, where we take 2 clean rows in which the shortcut disagrees with the label and pair each with a clean row of the same label in which the shortcut agrees, choosing the one closest in shortcut size $|x_8|$ (matched in size as well, so that the size does not give the label away either; \autoref{app:design}).

\paragraph{Drawing instead of ranking.}
In the exact targets, each shortcut-breaking row's size-matched partner scores as high as the row itself, so taking the top $k$ rows recovers balanced pairs. \sata cannot reproduce that pairing: it learns that shortcut-breaking rows deserve high scores, but not \emph{which} agreeing row is each one's partner, so it spreads the partners' share thinly over many agreeing rows. Taking its top $k$ rows therefore picks mostly shortcut-breaking ones, and the shortcut agreed with the label in only 27\% of the chosen rows: the shortcut was reversed again, with the stand-in learner's accuracy high under reversal (0.657) but at chance on ID queries (0.478). Because \sata is trained to output a distribution over rows (with a KL loss), it can instead \emph{draw} $k/2$ rows of each label in proportion to its scores; the spread-out partners then still make up about half of the draws (54\% agreement). On 50 validation tasks, drawing raised the worst environment on the stand-in learner from 0.478 to 0.558. Whether the drawing rule matters on Qwen itself is checked in \autoref{sec:sata} (\autoref{fig:rq4}, bottom).

\finding{A selector's training target can encode the very shortcut it is meant to avoid, and an oracle check on a stand-in learner can catch this before any training.}

\section{SATA on the model}
\label{sec:sata}

We use 1,000 tasks from the evaluation distribution to train \sata, regenerating them every epoch with a fresh permutation of feature positions so that the shortcut's position carries no information, and validate on 100 more (Adam \citep{kingma2015adam} at learning rate $10^{-4}$, early stopping \citep{prechelt1998early}).

On the stand-in learner its worst environment is 0.652, against 0.555 for the best hand-designed strategy, with no weak environment (\autoref{fig:sata-proxy}).

\finding{On the stand-in learner, \sata is the only method with no weak environment: it beats every hand-designed strategy under spurious reversal and stays within 0.03 of the best under ID and covariate shift.}

The \RQ4 family of six contrasts (\autoref{fig:rq4}; exact values in \autoref{tab:contrasts}) was decided and written down before any \sata result on the model.

\finding{On Qwen2.5-7B-Instruct, \sata does not beat balanced random demonstrations: none of the six pre-registered contrasts favours it.}

\begin{figure}[H]
\centering
\includegraphics[width=0.82\textwidth]{\figpath/fig_sata_proxy.pdf}
\caption{The stand-in learner's accuracy for each selection method, per environment (100 validation tasks), ordered by worst environment. Trained \sata matches its oracle targets and has no weak environment; the hand-designed strategies trade one environment against another.}
\label{fig:sata-proxy}
\end{figure}

\begin{figure}[H]
\centering
\includegraphics[width=\textwidth]{\figpath/fig_rq4.pdf}
\caption{\sata on Qwen2.5-7B-Instruct: the pre-registered \RQ4 family (top) and the check of drawing against top-$k$ selection (bottom), with 95\% intervals over 24 tasks. No contrast survives correction (all hollow); the interval for \sata against \strat{label\_diversity} under covariate shift lies below zero.}
\label{fig:rq4}
\end{figure}

\begin{table}[H]
\centering\small
\begin{tabular}{@{}lcccl@{}}
\toprule
\textbf{Strategy} & \textbf{ID} & \textbf{Covariate} & \textbf{Reversal} & \textbf{Shape} \\
\midrule
\strat{label\_diversity}  & 0.593 & 0.652 & 0.520 & good in two, weak under reversal \\
\sata                     & 0.553 & 0.556 & 0.552 & all three about equal \\
\strat{counter\_spurious} & 0.561 & 0.575 & 0.555 & about equal too \\
\bottomrule
\end{tabular}
\caption{Mean \AUROC on Qwen2.5-7B-Instruct under abstract names, by environment (24 tasks, three demonstration seeds).}
\label{tab:sata-flattening}
\end{table}

Based on the results above in \autoref{fig:rq4} and \autoref{tab:sata-flattening}, we can see that \sata tends to flatten the environments, bringing their performance to about the same level: it removes the weakness under reversal and pays for it in the ID and covariate environments. This is, however, expected, since \sata was trained to imitate the \strat{counter\_spurious} strategy, i.e.\ to make the shortcut uninformative. On ID and covariate queries, the shortcut is right 80 to 90\% of the time, and removing it loses that help, while under reversal it is wrong, so removing it avoids that harm.

The bottom group in \autoref{fig:rq4} compares \sata against itself, to see which of the two ways of turning its scores into 8 demonstrations works best. In the top-$k$ setting we chose the highest-scoring rows, 4 per label, whereas in the normal setting \sata draws rows in proportion to their scores. When we re-ran Qwen with \sata's top-$k$ rows instead, top-$k$ did no better than drawing, so the drawing rule is not why \sata did no better than balanced random selection on Qwen.


These results cannot fully separate a weakness of \sata from a property of the model, but they point towards the model: \sata's demonstrations carry the rule, since the stand-in learner learns it from them, yet Qwen does not use them. This is consistent with \citet{pan2023what} and \citet{wei2023larger}, who find that at the 7B scale in-context learning works mostly by recognising the task and applying what the model already believes, rather than by learning a new mapping from 8 examples. As our compute budget restricted us to 7B models, we aimed to see whether a less tuned model, with less pull from its priors, would give different results. Hence, in an exploratory study on Qwen2.5-7B (base), we ran the same \sata checkpoint and found that its weakest environment was better than any other strategy's, with its gains coming under spurious reversal (\autoref{app:base}). The base model also loses less to flipped names, i.e.\ it reads its demonstrations more. On a positive note, this means that with enough compute, \sata may perform differently, for better or worse, on larger models that can override their priors.

\finding{\sata helped the learners that read their demonstrations, the stand-in learner and, in an exploratory comparison, the base model, but not the instruction-tuned model.}

We leave readers with the possibility of training \sata with information about the underlying base model's behaviour, evaluated against that base model's priors. There may be room to improve \sata's demonstration selection if it were trained with this information in hand.

\chapter{Inside the model}
\label{ch:mechanism}

\emph{Research question: \RQ3, second half (why does the prior win?).}

\section{Motivation}

The two behavioural chapters converge on one picture: Qwen2.5-7B-Instruct does use its demonstrations, that is, any 8 labelled rows raise its \AUROC by about 0.1 over zero-shot under abstract and aligned names. But which of the 8 rows it sees barely mattered, whether the strategy was hand-designed or learned, while the feature names moved it by 0.2: aligned names lift every strategy together, as we saw in \autoref{ch:shift}, and flipped names, on the other hand, sink every strategy together. No choice of demonstrations seemed able to overturn the names. This raises the questions of whether the demonstrated rule is present in the model's internal state at all, whether a single direction in that state makes the model follow the demonstrations, at which layer the answer starts to depend on the names, and whether any of the selection strategies even move that layer.

\section{Method}
\paragraph{Logit lens.}

At every layer $l$ of 28, the last-position residual state $h_l$ (at the prompt's last token, the position from which the answer token is predicted) is passed through the model's final normalisation, and the difference of the two label rows of the unembedding is taken as $w = W_U[\texttt{1}] - W_U[\texttt{0}]$, giving a per-layer margin $m_l = w^\top \mathrm{norm}(h_l)$ \citep{nostalgebraist2020logit,elhage2021mathematical}. At the last layer, this is exactly the margin of \autoref{sec:prompts}.

\paragraph{Critical layer.}
Following \citet{halawi2024overthinking}, the critical layer is the first layer at which $\AUROC_l(\text{aligned}) - \AUROC_l(\text{flipped})$ exceeds 0.1 with a task-bootstrap interval above zero.
We define it with \AUROC rather than the accuracy of $\1[m_l > 0]$, because Qwen's bias toward label 1 would otherwise decide the accuracy (\autoref{sec:stats}).

\paragraph{Probes.}
A linear probe \citep{alain2017understanding} is a logistic regression that takes $h_l$ as input and predicts a label, where $h_l$ is the last-position state at layer $l$. If $h_l$ contains whatever Qwen took from the demonstrations, a probe should read the rule label from it; if it contains the answer Qwen would give without any demonstrations, a probe should read that instead. So we train two kinds of probe, rule and prior: across 15 layers (every second layer), 2 targets, 3 namings and 4 folds of training and testing, we come to 360 probes in total. To prevent the probes from memorising the training queries, we set a strong regulariser ($C = 0.1$).

\paragraph{Steering.}
The idea of activation steering was to see whether there was a way to switch something in Qwen so that it follows the demonstrations instead of its priors, because the logit lens and probes only help us observe what is happening in the model. Using contrastive activation addition \citep{rimsky2024steering}, we build $v_l$ as the difference of averages at layer $l$, computed for each rule label and then averaged over the two. On conflict queries, where the rule and Qwen's zero-shot guess disagree, the queries are split into two groups: those where Qwen's answer with demonstrations matched the rule, and those where it matched its no-demonstration guess instead. We build $v_l$ from the random-demonstration runs and re-run the queries under aligned and flipped names with the direction added at a chosen strength, in the manner of activation addition \citep{turner2023activation}. It is fitted on the 12 even-indexed tasks, which are the loan tasks, and applied, as $h_l \leftarrow h_l + \alpha v_l$ during the query only, on the 12 odd-indexed tasks, which are the medical ones, at layers 8, 12, 16, 20 and 24, with $\alpha$ from 1 to 32 and a random direction of the same norm as a control.

\section{Results}

\begin{figure}[H]
\centering
\includegraphics[width=\textwidth]{\figpath/fig_lens_strategies.pdf}
\caption{Logit-lens \AUROC against the rule label at every layer (24 tasks, seed 0), one line per strategy. The read-out stays near chance for 20 layers and splits by naming from layer 22 (dotted line: layer 23).}
\label{fig:lens}
\end{figure}

\finding{Qwen settles its answer only in its last few layers, and there the names, not the demonstrations, decide it.}

For the first 20 of Qwen's 28 layers the logit lens, which reads the answer Qwen would give if it stopped at that layer, is near chance (\AUROC 0.45--0.56) under every naming and strategy (\autoref{fig:lens}).
From layer 22 it splits by naming: with aligned names it rises to 0.63--0.67 (0.59 zero-shot), and with flipped names it falls to 0.45--0.52 (0.45 zero-shot).
The gap between the two passes 0.1 at layer 23 with random demonstrations, 24 zero-shot and 25 with \sata (\autoref{fig:lens-gap}).
Since the split appears without demonstrations too, the names cause it.
Demonstrations only change its size: random ones widen the final gap from 0.141 to 0.192, and the strategies built against the prior narrow it (\strat{counter\_prior} 0.174, \sata 0.104) without moving it earlier.

\begin{figure}[H]
\centering
\includegraphics[width=0.82\textwidth]{\figpath/fig_lens_gap.pdf}
\caption{Aligned minus flipped logit-lens \AUROC by layer, with a 95\% band over tasks for random demonstrations. The gap passes 0.1 at layer 23 with random demonstrations, 24 zero-shot and 25 with \sata; zero-shot shows the same jump, so the names drive it.}
\label{fig:lens-gap}
\end{figure}

\finding{Qwen's own guess can be read from its hidden states from layer 2 onward; the demonstrated rule hardly at all.}

Probes, simple classifiers trained to read one thing from the hidden states on some tasks and tested on others, read Qwen's zero-shot answer from layer 2 (\AUROC 0.80--0.87 with abstract names, 0.65--0.74 with named features), but the rule's label hardly at all (0.43--0.59; \autoref{fig:probes}).
A probe shows that information is present, not that it is used \citep{belinkov2022probing}, and a probe across tasks with different rules is a hard test; still, the contrast is stark: Qwen carries its own guess from its first layers, and the demonstrated rule barely registers.

\begin{figure}[H]
\centering
\includegraphics[width=\textwidth]{\figpath/fig_probes.pdf}
\caption{Linear probes on the hidden states (random demonstrations, tested on held-out tasks). Qwen's zero-shot answer is readable from layer 2; the rule label stays near chance.}
\label{fig:probes}
\end{figure}

\finding{Within the steering we could test (one averaged direction, five layers, a 7B model with eight demonstrations), pushing Qwen towards ``follow the demonstrations'' moves its answers without making them more correct.}

The steering direction is the average difference between Qwen's hidden states when it followed the demonstrations and when it followed its own guess; we added it, at strength $\alpha$, while Qwen read each query on 12 held-out tasks (the direction is small, 1--2\% of the hidden state's size, so the planned maximum of $\alpha = 4$ was raised to 32).
Up to $\alpha = 16$, \AUROC changes by at most 0.008; at $\alpha = 32$ a few results move by 0.02--0.04 either way, as they do for a random direction (\autoref{fig:steer}, top).
Yet at layer 24 the margins fall by about 9 logits (\autoref{fig:steer}, bottom): the push moves every answer the same way, so their order, which is all \AUROC measures, stays put.
It has to: following the demonstrations means lowering the margin when they imply 0 and raising it when they imply 1, and at layers 20 and 24 these two directions point almost opposite ways (cosine $-0.82$ to $-0.96$), so their average is a small leftover that shifts every answer together.

\begin{figure}[H]
\centering
\includegraphics[width=\textwidth]{\figpath/fig_steering.pdf}
\caption{Steering on the 12 held-out tasks: change in final \AUROC (top) and in the margin (bottom), with 95\% intervals over tasks. Blue shades are the ``follow the demonstrations'' direction at increasing strength $\alpha$; grey diamonds are a random direction of the same size. The margins fall by about 9 logits at layer 24 while \AUROC barely changes: the direction shifts every answer together.}
\label{fig:steer}
\end{figure}

\section{What this chapter does not show}

Neither the logit lens nor the probes is causal, and the lens can be biased in early layers \citep{belrose2023eliciting}, although the probes, which avoid that bias, agree with it.
Steering is the causal test, and it covers one kind of direction: a single averaged direction at five layers, added to the query.
It was fitted on the loan tasks and tested on the medical ones; refitting it on the medical tasks shows that this split does not explain the result.
\citet{minder2025controllable} find a context-versus-prior direction in Llama-3.1-8B; ours shows only that a direction built this way is not usable for this task and model, not that none exists.
A direction specific to each label, a learned one, steering of the demonstrations rather than the query, or a larger model could behave differently; the wider limitations are in \autoref{sec:limitations}.


\thesispart{Part III\quad Synthesis}
\chapter{Discussion}
\label{ch:discussion}

\section{Answers to the research questions}

\paragraph{\RQ1: does in-context learning degrade under shift?}
Yes, and concept shift, realised by naming, causes the most degradation. We also saw that reversing the shortcut costs about half as much, and unevenly: where the model's numeric prior agrees with the shortcut the cost is large, and where it opposes it there is almost none, because the prior reads the shortcut in the query with or without demonstrations. Averaged over tasks the prior's part cancels, and what remains is the shortcut the demonstrations taught; hiding it from the demonstrations shrinks the cost in every strategy and naming.

\begin{figure}[H]
\centering
\includegraphics[width=0.55\textwidth]{\figpath/fig_rq2_interaction.pdf}
\caption{The one \RQ2 contrast that survives Holm correction (abstract names). Each grey line is one task; blue is the mean with a 95\% interval (hierarchical bootstrap, 24 tasks). Neutralising the shortcut with \strat{counter\_spurious} costs \AUROC under covariate shift and roughly breaks even under its reversal; the rise between the two, $+0.11$ (Holm $p = 0.024$, 17 of 24 tasks), is the interaction.}
\label{fig:rq2-interaction}
\end{figure}

\newpage
\paragraph{\RQ2: do different strategies help under different shifts?}
Only weakly. The strategies differ far less than the namings do (\autoref{fig:strategies-by-naming}), and none lifts Qwen above chance under flipped names, including the one designed for prior conflict (+0.023 over \strat{label\_diversity}, not significant). Of the five planned \RQ2 contrasts, only one survives Holm correction (\autoref{fig:rq2-interaction}): \strat{counter\_spurious} fares better against \strat{label\_diversity} under shortcut reversal than under covariate shift, mainly because it costs \AUROC under covariate shift; its gain under reversal itself is not significant.

\paragraph{\RQ3: does the model rely on the right features, and why does the prior win?}

Under abstract names it tends to rely on the shortcut, whereas under aligned names it relies on the rule features, and under flipped names on the rule features but in the direction the names suggest. Its self-reports do not match the features it uses. The internal evidence fits the prior winning: the prior is present in the model's state from layer 2, while the demonstrated rule barely appears there, and the answer forms in the last six layers, where the names decide it (\autoref{ch:mechanism}).

\paragraph{\RQ4: does a learned selector help?}
Not on the instruction-tuned model, here Qwen2.5-7B-Instruct: none of the six pre-registered contrasts favours \sata, and taking its top-$k$ rows instead of drawing does no better. On the stand-in learner it is the only method with no weak environment, and in an exploratory comparison it helps the base model (\autoref{app:base}).

\section{The results in context}
\citet{wei2023larger} found that overriding semantic priors with in-context evidence emerges with scale. Our results are consistent with that picture and sharpen it for Qwen2.5-7B-Instruct and its base model. The instruction-tuned model is held more tightly by its priors (the base model loses less to flipped names, and a selector helps it). Although two models cannot show that instruction tuning is the cause, in what we tested the cost of conflict depended more on the names than on which evidence (demonstrations) was chosen. The graded picture of \citet{krishnakumar2025semantic} and \citet{kossen2024incontext}, in which mid-sized models follow flipped labels part of the time, fits the fact that demonstration labels still matter here (gold minus shuffled labels: +0.06 to +0.11 \AUROC, as seen in \autoref{ch:shift}) without being able to overturn the names.

\citet{halawi2024overthinking} find that under false demonstrations, the later layers tend to switch to the demonstrated labels. We find a late critical layer as well, but the switch it marks is driven mainly by the names, as the zero-shot runs show: in our setting the later layers are where the semantic priors are applied, not where the demonstrations take over. \citet{minder2025controllable} report a single context-versus-prior direction in Llama-3.1-8B but, for our numeric task, a mean-difference direction of that kind was a bias rather than a switch.

\citet{nastl2024causal} found that restricting predictors to causal features did not improve generalisation on real tabular shifts. We found the same for demonstrations without the shortcut on ID queries, although we also found a clear gain under reversal. Finally, learned selectors trained on model feedback \citep{rubin2022learning, ye2023compositional} are usually evaluated in distribution and on tasks the model already knows; our results, on the other hand, suggest that a selector's value depends on how much the model reads the evidence at all.

\section{Anticipated criticism}
\label{sec:criticism}

\subsection*{\RQ1}

% \begin{figure}[H]
% \centering
% \includegraphics[width=0.85\textwidth]{\figpath/.pdf}
% \label{fig:rq2-interaction}
% \end{figure}

\critique{Covariate shift showing no effect means your covariate shift is too weak.}

\critique{Naming is not concept shift; $P(y \mid x)$ never changes.}
{From the model's side it does: the same numbers under antonym names imply the opposite mapping, which is what concept shift means for a learner whose prior is semantic. We use the term in that operational sense and say so (\autoref{sec:environments}).}


\subsection*{\RQ2}

\critique{The strategies are weak straw men.}

Two observations suggest they are not. First, each follows a published idea (\autoref{tab:strategies}), and two of them, \strat{rule\_diversity} and \strat{counter\_spurious}, read the ground truth, which no deployed selector could. They are best-case versions of their ideas, so if they cannot beat balanced random, a realistic version is unlikely to, at least at $k = 8$. Second, the strategies are not inert: they change Qwen's behaviour the way their ideas predict. \strat{counter\_spurious}, which makes the shortcut uninformative, all but removes the drop from ID to spurious reversal under every naming: in \autoref{fig:strategies-by-naming} its ID and reversal dots nearly coincide in each panel (gaps of $-0.03$ to $0.01$, against $0.05$ to $0.07$ for \strat{label\_diversity}). It pays for this with lower ID \AUROC, so it does not beat \strat{label\_diversity} overall. The same holds for the others: no strategy reliably beats \strat{label\_diversity}, a random draw with $k/2$ rows of each label. Across six strategies, three namings and three environments, the typical difference from it is 0.025 \AUROC; 40 of the 54 differences have intervals that include zero, 10 lie below it and 4 above, and no strategy is clearly ahead in more than one environment under any naming. Each balanced strategy shows the same label counts as \strat{label\_diversity}, so these comparisons separate which rows are chosen from how many of each label are shown, which biases predictions on its own \citep{zhao2021calibrate}. Choosing rows by a selection idea therefore adds little over picking them at random with the labels balanced, and what it gains in one environment it tends to give back in another. The criticism does apply to selectors we did not test, such as retrievers trained on the model's own feedback; \RQ4 tests one learned selector.

\critique{24 tasks are too few to detect anything.}

They are too few to detect small effects, and we do not claim otherwise. The intervals for the strategy contrasts are about 0.1 \AUROC wide, so effects below about 0.05 would usually be missed; halving that width would take roughly four times as many tasks (assuming that the interval width shrinks with $\sqrt{n_T}$, where $n_T$ is the number of tasks). The design does detect effects of about 0.1: the cost of flipped names (0.20) and the \strat{counter\_spurious} interaction (0.11) both survive Holm correction, and \strat{feature\_range}'s loss under covariate shift ($-0.08$) has an interval entirely below zero (\autoref{tab:contrasts}). What the null results rule out is a large benefit. For the four strategies tested against \strat{label\_diversity} under their target shifts, the largest gain inside any interval is 0.09 \AUROC, under half the cost of flipped names, and two of the four exclude any gain above 0.05; the six \sata contrasts give the same ceiling (\autoref{ch:sata}). A strategy that helps by a few hundredths of \AUROC may exist at $k = 8$ on this model; one that helps as much as the names hurt is ruled out though.

\newpage

\subsection*{\RQ3}

\critique{The logit lens is unreliable for early layers.}

There is a possibility that it can be biased for early layers \citep{belrose2023eliciting}.
The final layer matches the model's output exactly, and the linear probes avoid that bias. The probes tell the same story: the prior is active early in the model, and the rule stays weak throughout.

\critique{Steering failing proves nothing about steering in general.}

Agreed. We tested one family of directions (a label-averaged mean difference at five layers, applied to the query) on one 7B model with eight numeric demonstrations, within a small budget. We claim only that this direction acts as a bias at this scale, because steering adds the same vector $v_l$ to every query, whatever its label, and we explain why it must: following the demonstrations means opposite things for the two labels. When the demonstrations imply label 0, following them means pushing the margin down; when they imply label 1, it means pushing it up, and the measured cosine between the two per-label directions was approximately $-0.9$, i.e.\ they point in nearly opposite directions. Within-cell \AUROC (a ``cell'' is our term for one fixed prompt setup: one task, one strategy, one demonstration seed, one naming and one environment) cannot see a push that moves every margin together, which is why \autoref{fig:steer} reports the margins as well: at layer 24 they fall by about 9 logits, towards answering 0, while \AUROC stays flat, so the direction shifts every answer at once rather than separating the labels.

\critique{Self-reported importance is a known failure; why test it?}

Because it is the cheapest explanation a practitioner would ask for, and because the comparison with behavioural reliance under each naming shows that the self-report follows the names, not the model's behaviour.

\newpage
\subsection*{\RQ4}
\critique{\sata failed because it was badly trained.}

It reached its own oracle's level on the stand-in learner and was the only method with no weak environment there; the base model benefits from the same checkpoint; and ranking instead of drawing does no better. The results point instead to what the instruction-tuned model does with the evidence it is given, as discussed earlier.

\critique{The stand-in learner is a poor proxy for the model, so the selector was optimised for the wrong thing.}

This is the main weakness of the selector chapter, and we measured it. On 8 held-out tasks shown to Qwen, the stand-in learner ranks the strategies across environments much as Qwen does under abstract names: we saw the correlation to be about 0.7, but only weakly under named features, about 0.5. Two results suggest that the proxy still targeted something real. \sata's demonstrations carry the rule, because the stand-in learner learns from them in every environment, and the base model also uses them: with the same checkpoint, its weakest environment is better than any other strategy's (exploratory; \autoref{app:base}). What the proxy did miss, however, is how little the instruction-tuned model uses the evidence once the names and its prior pull the other way. So the direct fix may be to train the selector on Qwen's own answers rather than on a stand-in: running Qwen with many random demonstration sets and scoring each row by how much it helps, just as \citet{chang2023data} do to find helpful demonstrations. However, the inherent bias of a model of this scale would also affect the demonstrations chosen by such a selector, so a larger model, whether in the Qwen family or another, would be more appropriate. Although this was planned, it could not be run within our budget.


\newpage
\section{Limitations}
\label{sec:limitations}

Every negative result in this thesis holds for the conditions we could test, and those conditions were narrow.

\begin{itemize}[leftmargin=1.5em]
  \item \textbf{Scale.} One model family at 7B: Qwen2.5-7B-Instruct and, exploratorily, its base model. Due to time and budget limitations, Llama-3.1-8B-Instruct was dropped when it would not answer in format, and larger models, which \citet{wei2023larger} find follow contradicting evidence better, were out of reach.
  \item \textbf{Demonstrations.} $k = 8$, chosen on the pilot because $k = 32$ did not pass the learnability check. More demonstrations might give the rule more weight.
  \item \textbf{Tasks.} Synthetic numeric tasks with three load-bearing features, two domains and 24 evaluation tasks. Real tabular data, text features or other rule families may behave differently.
\item \textbf{Metric.} Within-cell \AUROC cannot see a shift that moves every query's margin by the same amount; balanced accuracy, its alternative, was distorted by calibration.
  \item \textbf{Proxies.} The selector was trained against a stand-in learner that tracks the model only moderately under named features.
\item \textbf{Budget.} The whole project used a few tens of US dollars of rented GPU time, and the later stages were sized to fit. Steering tested one kind of direction (a single averaged direction, added to the query). Two selectors that learn from Qwen's own answers rather than from a stand-in were planned but not run, because each needs many runs of Qwen for every training task: one that scores each candidate row by how much it helps Qwen \citep{chang2023data}, and one trained by reinforcement learning to add demonstrations one at a time \citep{zhang2022active}. They are the most direct test of whether the stand-in learner, rather than the model, held \sata back.

\item \textbf{Decisions after data.} Some design choices were made after seeing early data. Most never touched the evaluation tasks: the shortcut's direction was randomised per task after the pilot showed Qwen's ``bigger means 1'' habit backing it; the feature names were picked by their measured effect on Qwen's zero-shot answers on random profiles; and \sata's matched target and drawing rule were chosen on its own validation tasks, judged by the stand-in learner. However, the ten planned tests for \RQ1--\RQ3 had been written down in advance with calibrated balanced accuracy as their measure, and were later switched to \AUROC because, after calibration, a large share of the prompts (about 36\%) gave the same answer to every query.
\end{itemize}

\section{Future work}
The most immediate direct extension is scale, with the same testbed and contrast families on larger models, where the influence of the priors should be weaker. This would show whether selection starts to matter as the model starts to read the evidence. Beyond this, datamodels \citep{ilyas2022datamodels,chang2023data}, where we would train the selector on the model's own scores of demonstration subsets, possibly read at the critical layer, would remove the proxy problem described earlier. One of the parked ideas is to turn the prior rather than fight it, by negating the features whose larger values go with a specific label, in demonstrations and query alike.


\chapter{Conclusion}
\label{ch:conclusion}

The thesis asked whether choosing demonstrations can help a relatively small model, at 7B, follow them over its own priors under controlled distribution shift. To answer it, we built a synthetic tabular testbed in which covariate shift, spurious reversal and concept shift can each be switched on alone, set the model's semantic prior against the rule by naming the features, and tested every confirmatory claim against a contrast family frozen before its data were seen.

For Qwen2.5-7B-Instruct with 8 numeric demonstrations, none of the selection methods we tried did so. The feature names decided most of what the model predicted, more than any shift of the data. The hand-designed strategies differed little from balanced random demonstrations; the one clear effect is that neutralising the shortcut fares better under its reversal than under covariate shift. A learned selector, the only method with no weak environment on a stand-in learner, did no better than balanced random demonstrations on the model, though in an exploratory study it did help the base model, whose priors pull less. Inside the model, the answers form in the last few layers, and the names decide them there. The prior is present from the start and the demonstrated rule hardly at all, and a single steering direction moves the answers without making them more correct.

Selection chooses evidence; in our experiments, it did not make the instruction-tuned model read evidence it largely ignored. Within the scale we could study, this appears to be the binding constraint. The testbed we have constructed here, including the pre-registered families and the diagnosis, gives future work, at a larger scale and with the model's own training signal, a controlled place to find out when it stops being a constraint.


%%%%%%%%%%%%
% End

% Bibliography
{
\setstretch{1.0}\small
\cleardoublepage
\phantomsection
\printbibliography[heading=bibintoc,title={References}]
}

%%%% Appendices
\appendix
\chapter{Notation}
\label{app:notation}

Symbols are grouped by where they first appear.
Every symbol is also defined in the text at its first use; this appendix is the single place to look one up.
Two letters carry more than one meaning, always told apart by a subscript or superscript:
$\sigma(\cdot)$ is the logistic function and $\sigma_t$ a task's spurious direction;
$s$ is the shortcut's strength, $s_j$ a load-bearing feature's direction in the rule, and $s^{\mathrm{lex}}$ a feature name's dictionary direction.

{\small
\renewcommand{\arraystretch}{1.05}
\begin{longtable}{@{}p{3.4cm}p{11.6cm}@{}}
\toprule
\textbf{Symbol} & \textbf{Meaning} \\
\midrule
\endhead
\multicolumn{2}{@{}l}{\textit{Tasks and data}} \\
$f_0, \ldots, f_9$ & The ten features of a task, in the generator's numbering; also their abstract display names \\
$x \in \mathbb{R}^{10}$, $x_j$ & A row's feature vector and its $j$-th feature, standardised with the pool's statistics \\
$C$, $|C| = 3$ & The load-bearing (causal) features: the label depends on these and nothing else \\
$x_C$ & The load-bearing part of $x$ \\
$c_j = u_j r_j$ & Linear coefficient of load-bearing feature $j$, with $u_j \sim \mathrm{U}(1,2)$ and a random sign $r_j \in \{+1,-1\}$ \\
$s_j \in \{+1,-1\}$ & Direction of load-bearing feature $j$ in the rule ($s_j = \operatorname{sign} c_j$ for linear tasks) \\
$\sigma(z) = 1/(1+e^{-z})$ & The logistic function \\
$\mathrm{Maj}$ & Majority vote over three signs (tree rules) \\
$y_{\mathrm{clean}}$, $y$ & A row's rule label before and after label noise \\
$\eta = 0.02$ & Label-noise rate: the probability that $y \neq y_{\mathrm{clean}}$ \\
$s \sim \mathrm{U}(0.80, 0.90)$ & Spurious strength: the probability that the shortcut's sign agrees with the label \\
$d = \Phi^{-1}(s)$ & The shortcut's class-conditional mean, with $\Phi$ the standard normal distribution function \\
$\sigma_t \in \{+1,-1\}$ & Task $t$'s spurious direction: whether a large shortcut value indicates label 1 or label 0 \\
$\varepsilon$ & Standard normal noise in the shortcut \\
$\pi_t$ & Task $t$'s permutation of feature positions (the displayed order) \\
$T$, $n$ & Number of tasks; queries per task and environment ($T = 24$, $n = 20$ in evaluation) \\
$\delta$ & Size of a shift or a nudge, in standard deviations \\
\multicolumn{2}{@{}l}{\textit{Naming, prompts and the model}} \\
$N$ & A naming: abstract, aligned or flipped \\
$s^{\mathrm{lex}} \in \{+1, 0, -1\}$ & A display name's dictionary direction ($+1$ for the first member of an antonym pair, $-1$ for the second, 0 for neutral and abstract names) \\
$k$ & Demonstrations per prompt ($k = 8$) \\
$\ell_0$, $\ell_1$ & The model's logits for the label tokens \texttt{0} and \texttt{1} \\
$m = \ell_1 - \ell_0$ & The label margin, equal to $\log p(\texttt{1}) - \log p(\texttt{0})$ \\
$m_0(x; N)$ & The model's zero-shot margin on row $x$ under naming $N$: its prior \\
$\tilde m_0(x_i; N)$ & The measured prior of pool row $i$, centred on the pool median \\
$C_t$ & Task $t$'s prior--data conflict under naming $N$: the share of pool rows whose label disagrees with the prediction of the centred prior, $\1[\tilde m_0 > 0]$ \\
% \multicolumn{2}{@{}l}{\textit{Evaluation and statistics}} & mentioned it as is by definition \\

\textit{Evaluation and statistics} & mentioned it as is by definition \\

\AUROC & Area under the ROC curve of the margin against the label, within one cell (task, strategy, seed, naming, environment) \\
$B = 2000$ & Bootstrap draws \\
$\pi_{\mathrm{true}}$, $\pi_{\mathrm{behav}}$, $\pi_{\mathrm{self}}$ & Feature importance: ground truth, behavioural (margin drop when the feature is replaced), and the model's self-report \\
$\rho(\cdot,\cdot)$ & Spearman rank correlation \\
$g_j$ & Directional sensitivity: $m(x + \delta e_j) - m(x - \delta e_j)$ with $\delta = 0.5$ \\
DFI & Demo-following index: the share of (query, load-bearing feature) pairs with $\operatorname{sign} g_j = s_j$ \\
\multicolumn{2}{@{}l}{\textit{The selector}} \\
$p_\theta$ & The selector's softmax distribution over pool rows, with parameters $\theta$ \\
$q$ & Its training target: a distribution over pool rows \\
$\mathrm{KL}(q \,\Vert\, p_\theta)$ & The training loss \\
\multicolumn{2}{@{}l}{\textit{Inside the model}} \\
$L = 28$, $l$ & Number of decoder layers; a layer index ($l = 0$ is the embeddings) \\
$h_l$ & The last-position residual state after layer $l$ \\
$W_U$ & The unembedding (output) matrix \\
$w = W_U[\texttt{1}] - W_U[\texttt{0}]$ & The label direction: the difference of the two label rows \\
$\mathrm{RMSNorm}$ & The model's final normalisation layer \\
$m_l$ & The logit-lens margin at layer $l$: $w^\top \mathrm{RMSNorm}(h_l)$ for $l < L$, and $m_L = m$ \\
$\AUROC_l$ & \AUROC of $m_l$ against the rule label \\
$v_l$ & The steering vector at layer $l$ \\
$\alpha$ & Steering strength: $h_l \leftarrow h_l + \alpha v_l$ \\
\bottomrule
\end{longtable}
}
\chapter{Design details}
\label{app:design}

This appendix gives the definitions and properties that the chapters rely on.

\section{Tasks and environments}

\paragraph{Rows.}
Features are drawn as $x \sim \mathcal{N}(0, I_{10})$, the environment's perturbation is applied, the clean label is computed from the load-bearing features $C$, the shortcut $x_8$ is generated from the clean label, $x_9 \sim \mathcal{N}(0,1)$, and labels are flipped with probability $\eta = 0.02$.
Features are standardised with the pool's statistics and shown in a random order $\pi_t$ per task, so no role can be read from a position.

\paragraph{Rules.}
A \emph{linear} task has $y_{\mathrm{clean}} \sim \mathrm{Bernoulli}\big(\sigma(c^\top x_C)\big)$ with $c_j = u_j r_j$, $u_j \sim \mathrm{U}(1,2)$, $r_j \in \{\pm 1\}$; $P(y_{\mathrm{clean}} = 1 \mid x)$ is strictly monotone in each $x_j$, so feature $j$ has the direction $s_j = \operatorname{sign} c_j$.
A \emph{tree} task splits each of three features at zero and labels the eight leaves; the leaf tables allowed (complement-antisymmetric, neither a dictator nor parity) are exactly the signed majorities $y_{\mathrm{clean}} = \mathbb{1}\big[\mathrm{Maj}(s_1 x_1, s_2 x_2, s_3 x_3) > 0\big]$, so tree tasks are monotone too.

\paragraph{The shortcut.}
With $d = \Phi^{-1}(s)$, $x_8 = \sigma_t\big((2y_{\mathrm{clean}} - 1)d + \varepsilon\big)$, $\varepsilon \sim \mathcal{N}(0,1)$.

\begin{proposition}[Sign agreement]
$P\big(\operatorname{sign}(\sigma_t x_8) = \operatorname{sign}(2y_{\mathrm{clean}}-1)\big) = s$; under spurious reversal ($d \to -d$) it is $1 - s$.
\end{proposition}
\begin{proof}
For $y = 1$, $P(\sigma_t x_8 > 0) = P(Z > -d) = \Phi(d) = s$; $y = 0$ is symmetric. Reversal replaces $d$ by $\Phi^{-1}(1-s) = -d$.
\end{proof}

\paragraph{Covariate shift.}
Two load-bearing features $i, j$ are moved in opposite label directions and a distractor $d \notin C$ at random:
$\delta_i = s_i m/|c_i|$, $\delta_j = -s_j m/|c_j|$, $\delta_d \sim \pm\mathrm{U}(1,2)$, with the push $m$ chosen so that every shifted feature moves by 1--2 standard deviations (for trees $|c| = 1$, $m \sim \mathrm{U}(1,2)$).

\begin{lemma}[The label rate is unchanged]
(a) For linear tasks, $c^\top(x + \delta) = c^\top x$ for every row. (b) For tree tasks, $P(y_{\mathrm{clean}} = 1)$ stays $1/2$.
\end{lemma}
\begin{proof}
(a) $c_i\delta_i = m$, $c_j\delta_j = -m$, and the distractor carries no weight. (b) The three votes are independent. After the shift the two moved votes are 1 with probabilities $p_i = \Phi(m)$ and $p_j = 1 - \Phi(m)$, and the third with $\tfrac12$, so the majority is 1 with probability $p_i p_j + \tfrac12[p_i(1-p_j) + (1-p_i)p_j] = \tfrac12(p_i + p_j) = \tfrac12$.
\end{proof}

\section{Naming}

Each domain has antonym pairs $(p, q)$, where $p$ is associated in everyday language with label 1 and $q$ with label 0, and a set of weak-prior names.
Each load-bearing feature $j$ is assigned a pair: \emph{aligned} names it $p$ if $s_j = +1$ and $q$ if $s_j = -1$; \emph{flipped} uses the other member.
Every other feature gets a weak-prior name, the same in both conditions.

\begin{proposition}[Identification]
For a fixed task, demonstration set and query, the aligned and flipped prompts differ only in the names of the three load-bearing features. A predictor that ignores names has the same expected accuracy under both, so a difference in the model's accuracy is attributable to the meaning of the names.
\end{proposition}

\section{Counter-spurious selection}

Balancing the four (label, shortcut agrees) cells is not enough: agreeing rows lie further from zero ($\mathbb{E}[|x_8| \mid \text{agree}] = d + \phi(d)/\Phi(d) > \mathbb{E}[|x_8| \mid \text{disagree}] = \phi(d)/(1-\Phi(d)) - d$), so $x_8$ still correlates with the label (about 0.33 at $s = 0.85$).
\strat{counter\_spurious} therefore takes, within each label, $k/4$ clean disagreeing rows and pairs each with the clean agreeing row of the same label nearest in $|x_8|$ \citep{rubin1973matching,stuart2010matching}.
With perfect matching, sign agreement is $\tfrac12$ and $\mathrm{Cov}(x_8, 2y-1) = 0$; on the evaluation tasks the implementation gives $\mathrm{corr}(x_8, y) = 0.005$.
\sata's matched target (\autoref{sec:sata-targets}) is each row's probability $\pi_i$ of appearing in such a set built for the query, with rows in the query's regime (the sign pattern of its three load-bearing features) taken first, normalised: $q_i = \pi_i / \sum_j \pi_j$; label-noise rows have $\pi_i = 0$, and nothing reads the query's label.
\sata is trained with the loss $\mathrm{KL}(q \,\Vert\, p_\theta)$.

\section{Scoring and caching}

\begin{lemma}[Margin]
With $p_1 = e^{\ell_1}/(e^{\ell_0} + e^{\ell_1})$ and $m = \ell_1 - \ell_0$, $p_1 = \sigma(m)$, and $m$ is the same whether computed from logits or from full-vocabulary log-probabilities.
\end{lemma}
\begin{proof}
$p_1/(1 - p_1) = e^{\ell_1 - \ell_0}$; log-probabilities subtract a common log-sum-exp, which cancels.
\end{proof}

\begin{proposition}[Prefix caching is exact]
For a decoder-only transformer with causal attention, running the prefix $P$ once and each suffix $S$ against its cached keys and values gives the same final-position logits as a full pass on $P \Vert S$, provided $\mathrm{tok}(P \Vert S) = \mathrm{tok}(P) \Vert \mathrm{tok}(S)$.
\end{proposition}
\begin{proof}
States at positions $< |P|$ depend only on those positions' tokens (attention reads positions $\le t$), so appending $S$ does not change them, and the states of $S$ are computed from the same inputs either way. The runner checks the tokenisation condition for every suffix.
\end{proof}

\begin{lemma}[Logit lens at the last layer]
With $w = W_U[\texttt{1}] - W_U[\texttt{0}]$ and $h_L$ the final (normalised) state, $m_L = w^\top h_L = \ell_1 - \ell_0 = m$.
\end{lemma}

\section{Statistics}

\paragraph{Cells and contrasts.}
\AUROC is computed within each cell (task, strategy, seed, naming, environment) over its 20 label-balanced queries and averaged over seeds.
A contrast is a per-task difference of such values (for example, strategy $A$ minus strategy $B$ on the same queries, or ID minus a shifted environment), averaged over the 24 tasks.

\paragraph{Hierarchical bootstrap.}
Each of $B = 2000$ draws resamples tasks with replacement and, within each resampled task, its positive and negative queries separately; the percentile interval of the resampled estimates is reported, with one-sided $p$-values from the share of draws on the wrong side of zero.
Because tasks differ in their rules, the variance of the mean is dominated by the between-task term ($\sigma_T^2/T + \sigma_W^2/(Tn)$; \citealp{cochran1977sampling}), so tasks are the unit of resampling.

\paragraph{Multiplicity.}
Holm's step-down correction \citep{holm1979simple} is applied within each pre-registered family: ten contrasts for \RQ1--\RQ3 and six for \sata.
\autoref{tab:contrasts} lists every confirmatory estimate.

\begin{table}[h]
\centering\small
\caption{Every confirmatory contrast (Qwen2.5-7B-Instruct, 24 tasks; one-sided hierarchical bootstrap, Holm within each family). All are differences in \AUROC except demo-following (in the demo-following index) and feature reliance (in Spearman $\rho$). Bold: significant after correction.}
\label{tab:contrasts}
\begin{tabular}{@{}lrrr@{}}
\toprule
\textbf{Contrast} & \textbf{Estimate [95\% CI]} & \textbf{Holm $p$} & \textbf{Tasks} \\
\midrule
\multicolumn{4}{@{}l}{\emph{Shift, names and the prior (\RQ1--\RQ3)}} \\
Covariate shift: ID $-$ covariate (random, abstract) & $+0.024$ [$-0.044$, 0.093] & 0.95 & 15/24 \\
Shortcut reversal: ID $-$ reversal (random, abstract) & $+0.115$ [0.017, 0.218] & 0.063 & 16/24 \\
\textbf{Flipped names: aligned $-$ flipped (random, ID)} & $\mathbf{+0.204}$ [0.131, 0.280] & \textbf{0.005} & 22/24 \\
counter\_spurious $-$ label\_diversity, reversal & $+0.035$ [$-0.019$, 0.086] & 0.59 & 16/24 \\
similarity $-$ label\_diversity, covariate & $-0.029$ [$-0.110$, 0.048] & 1.00 & 11/24 \\
feature\_range $-$ label\_diversity, covariate & $-0.078$ [$-0.141$, $-0.018$] & 1.00 & 6/24 \\
counter\_prior $-$ label\_diversity, flipped, ID & $+0.023$ [$-0.015$, 0.064] & 0.70 & 7/24 \\
\textbf{counter\_spurious's gain: reversal $-$ covariate} & $\mathbf{+0.111}$ [0.033, 0.196] & \textbf{0.024} & 17/24 \\
\textbf{Demo-following index: abstract $-$ flipped} & $\mathbf{+0.123}$ [0.079, 0.175] & \textbf{0.005} & 23/24 \\
Feature reliance $\rho(\pi_{\mathrm{true}}, \pi_{\mathrm{behav}})$: abstract $-$ flipped & $+0.026$ [$-0.141$, 0.174] & 1.00 & 13/24 \\
\midrule
\multicolumn{4}{@{}l}{\emph{\sata (\RQ4)}} \\
\sata{} $-$ label\_diversity, reversal & $+0.032$ [$-0.024$, 0.089] & 0.79 & 14/24 \\
\sata{} $-$ label\_diversity, covariate & $-0.096$ [$-0.154$, $-0.041$] & 1.00 & 5/24 \\
\sata{} $-$ label\_diversity, ID & $-0.040$ [$-0.102$, 0.023] & 1.00 & 8/24 \\
\sata{} $-$ similarity, mean of 3 environments & $-0.020$ [$-0.064$, 0.026] & 1.00 & 11/24 \\
\sata{} $-$ \sata{} without the query, mean of 3 & $-0.025$ [$-0.063$, 0.010] & 1.00 & 10/24 \\
\sata{} $-$ label\_diversity, flipped, ID & $-0.014$ [$-0.077$, 0.048] & 1.00 & 13/24 \\
\bottomrule
\end{tabular}
\end{table}

\section{Reproducibility}
\label{app:repro}

The code, configuration and pre-registered contrast families accompany this thesis, with tests covering the invariants of the generator, the selection engine, the model runner and the probes.
Every prediction is stored with its demonstration rows, order and prompt, so every number and figure here can be recomputed from the result files.\todo{Add the repository link.}
\chapter{Exploratory results: \sata on the base model}
\label{app:base}

This appendix reports the exploratory comparison summarised in \autoref{sec:sata}.
It was not part of a pre-registered family, so it is a lead for future work rather than a confirmed result.

\paragraph{Setup.}
We ran the same \sata checkpoint on Qwen2.5-7B (base), on the 24 evaluation tasks, as part of the reduced base-model grid of \autoref{ch:shift}: abstract and flipped names only, and three of the hand-designed strategies (\strat{random}, \strat{counter\_spurious} and \strat{counter\_prior}) besides zero-shot.
\strat{label\_diversity} was not run on the base model, so \sata is compared here with \strat{random} and \strat{counter\_spurious}.

\begin{figure}[H]
\centering
\includegraphics[width=0.8\textwidth]{\figpath/fig_base_sata.pdf}
\caption{Every strategy run on Qwen2.5-7B (base), by naming (panels) and environment (dot colour), as mean \AUROC over 24 tasks. The grey line joins one strategy's three environments, and rows are sorted by each strategy's weakest environment under abstract names. \sata's three environments sit close together and its weakest one is the highest of any strategy; the hand-designed strategies trade one environment against another.}
\label{fig:base-sata}
\end{figure}

\paragraph{What the figure shows.}
Under both namings, \sata's three environments sit close together, and its weakest environment is higher than any other strategy's (\autoref{fig:base-sata}).
The hand-designed strategies trade one environment against another: \strat{random} and \strat{counter\_prior} do well on ID and covariate queries and fall under spurious reversal, the same pattern as on the instruction-tuned model (\autoref{tab:sata-flattening}), while \strat{counter\_spurious} is flatter but lower.
This is the flattening \sata showed on the instruction-tuned model, but on the base model it comes without a clear loss on ID and covariate queries.

\paragraph{The paired comparisons.}
\autoref{tab:base-contrasts} compares \sata with the two strategies on the same queries.
Against \strat{random}, its gains come under spurious reversal, with no clear cost on ID or covariate queries.
Against \strat{counter\_spurious}, it is ahead in all six cells, with intervals above zero in two.
Only one comparison, the gain over \strat{random} under spurious reversal with flipped names, survives a Holm correction over the twelve.

\begin{table}[H]
\centering\small
\caption{\sata against the two hand-designed strategies run on the base model (24 tasks; hierarchical bootstrap as in \autoref{sec:stats}; Holm over these twelve comparisons). Exploratory. Tasks: the number of tasks in which \sata is ahead. Bold: survives correction.}
\label{tab:base-contrasts}
\begin{tabular}{@{}llrrr@{}}
\toprule
\textbf{Names} & \textbf{Comparison, environment} & \textbf{Estimate [95\% CI]} & \textbf{Holm $p$} & \textbf{Tasks} \\
\midrule
abstract & \sata{} $-$ random, ID & $-0.027$ [$-0.095$, 0.040] & 1.00 & 13/24 \\
 & \sata{} $-$ random, covariate & $-0.014$ [$-0.080$, 0.051] & 1.00 & 12/24 \\
 & \sata{} $-$ random, reversal & $+0.085$ [0.015, 0.155] & 0.11 & 15/24 \\
 & \sata{} $-$ counter\_spurious, ID & $+0.048$ [$-0.017$, 0.110] & 0.45 & 17/24 \\
 & \sata{} $-$ counter\_spurious, covariate & $+0.050$ [$-0.011$, 0.108] & 0.45 & 14/24 \\
 & \sata{} $-$ counter\_spurious, reversal & $+0.029$ [$-0.028$, 0.087] & 0.94 & 14/24 \\
\midrule
flipped & \sata{} $-$ random, ID & $+0.008$ [$-0.060$, 0.077] & 1.00 & 16/24 \\
 & \sata{} $-$ random, covariate & $0.000$ [$-0.072$, 0.076] & 1.00 & 12/24 \\
 & \textbf{\sata{} $-$ random, reversal} & $\mathbf{+0.094}$ [0.031, 0.159] & \textbf{0.024} & 19/24 \\
 & \sata{} $-$ counter\_spurious, ID & $+0.077$ [0.012, 0.136] & 0.12 & 19/24 \\
 & \sata{} $-$ counter\_spurious, covariate & $+0.079$ [0.003, 0.152] & 0.18 & 18/24 \\
 & \sata{} $-$ counter\_spurious, reversal & $+0.020$ [$-0.038$, 0.077] & 1.00 & 15/24 \\
\bottomrule
\end{tabular}
\end{table}

\paragraph{Limits.}
These comparisons were not pre-registered; the base grid covers two namings and three of the seven hand-designed strategies; and the comparison is between two models only, so it cannot show that instruction tuning is what makes the difference.
The base model also loses less to flipped names than the instruction-tuned model (\autoref{ch:shift}), which fits the reading that \sata helps a model only as far as it reads its demonstrations, but does not establish it.


\end{document}
